% Proyecto de Grado — Pontificia Universidad Javeriana Cali.
% Compilar: latexmk -pdf main.tex
% Estructura de escritura: completar los marcadores rojos siguiendo los comentarios GUÍA.
\documentclass[11pt,oneside]{thesis}
\usepackage{enumerate}
\usepackage{longtable}
\usepackage{url}
\usepackage{enumitem}
\usepackage{float}
\usepackage{booktabs}
\usepackage{listings}
\usepackage{caption}
\usepackage{subcaption}
\usepackage{makeidx}
\usepackage{graphicx}
\usepackage{tabularx}
\usepackage{pgfgantt}
\usepackage[spanish]{babel}
\usepackage[utf8]{inputenc}
\usepackage{hyperref}

% GUÍA: Datos de portada y cartas tomados de las instrucciones del autor.
\titulo{Sistema de apoyo al diagnóstico de melanoma mediante aprendizaje profundo e integración en aplicativo web}
\autorA{\pendiente{código}}{Julián Andrés Galvis Tocora}
\fecha{Noviembre de 2026}
\directorproyecto{PhD.}{Hernán Darío Vargas Cardona}
\directorcarrera{Dr.}{Gerardo Mauricio Sarria} % VERIFICAR nombre vigente

\usepackage{amsmath,amssymb}             % AMS Math
% \usepackage[french]{babel}
%\usepackage[latin1]{inputenc}
%\usepackage[T1]{fontenc}
%\usepackage[utf8]{inputenc}
\usepackage[left=1.0in,right=1.0in,top=1.0in,bottom=1.0in,includefoot,includehead,headheight=13.6pt]{geometry}
\renewcommand{\baselinestretch}{1.05}

% Table of contents for each chapter

\usepackage[nottoc, notlof, notlot]{tocbibind}
\usepackage{minitoc}
\setcounter{minitocdepth}{2}
\mtcindent=15pt
% Use \minitoc where to put a table of contents

\usepackage{aecompl}

% My pdf code

\usepackage{ifpdf}


% Links in pdf
\usepackage{color}
\definecolor{linkcol}{rgb}{0,0,0.4} 
\definecolor{citecol}{rgb}{0.5,0,0} 

% Change this to change the informations included in the pdf file

% See hyperref documentation for information on those parameters

\hypersetup
{
bookmarksopen=true,
pdftitle={\titulovar},
pdfauthor={\authorA},
pdfsubject={Proyecto de Grado},
%pdftoolbar=false, % toolbar hidden
pdfmenubar=true, %menubar shown
pdfhighlight=/O, %effect of clicking on a link
colorlinks=true, %couleurs sur les liens hypertextes
pdfpagemode=UseNone,
pdfpagelayout=SinglePage, %ouverture en simple page
pdffitwindow=true, %pages ouvertes entierement dans toute la fenetre
linkcolor=linkcol, %couleur des liens hypertextes internes
citecolor=citecol, %couleur des liens pour les citations
urlcolor=linkcol %couleur des liens pour les url
}

% definitions.
% -------------------

\setcounter{secnumdepth}{3}
\setcounter{tocdepth}{2}

% Some useful commands and shortcut for maths:  partial derivative and stuff

\newcommand{\pd}[2]{\frac{\partial #1}{\partial #2}}
\def\abs{\operatorname{abs}}
\def\argmax{\operatornamewithlimits{arg\,max}}
\def\argmin{\operatornamewithlimits{arg\,min}}
\def\diag{\operatorname{Diag}}
\newcommand{\eqRef}[1]{(\ref{#1})}

\usepackage{rotating}                    % Sideways of figures & tables
%\usepackage{bibunits}
%\usepackage[sectionbib]{chapterbib}          % Cross-reference package (Natural BiB)
%\usepackage{natbib}                  % Put References at the end of each chapter
                                         % Do not put 'sectionbib' option here.
                                         % Sectionbib option in 'natbib' will do.
\usepackage{fancyhdr}                    % Fancy Header and Footer

% \usepackage{txfonts}                     % Public Times New Roman text & math font
  
%%% Fancy Header %%%%%%%%%%%%%%%%%%%%%%%%%%%%%%%%%%%%%%%%%%%%%%%%%%%%%%%%%%%%%%%%%%
% Fancy Header Style Options

\pagestyle{fancy}                       % Sets fancy header and footer
\fancyfoot{}                            % Delete current footer settings

%\renewcommand{\chaptermark}[1]{         % Lower Case Chapter marker style
%  \markboth{\chaptername\ \thechapter.\ #1}}{}} %

%\renewcommand{\sectionmark}[1]{         % Lower case Section marker style
%  \markright{\thesection.\ #1}}         %

\fancyhead[LE,RO]{\bfseries\thepage}    % Page number (boldface) in left on even
% pages and right on odd pages
\fancyhead[RE]{\bfseries\nouppercase{\leftmark}}      % Chapter in the right on even pages
\fancyhead[LO]{\bfseries\nouppercase{\rightmark}}     % Section in the left on odd pages

\let\headruleORIG\headrule
\renewcommand{\headrule}{\color{black} \headruleORIG}
\renewcommand{\headrulewidth}{1.0pt}
\usepackage{colortbl}
\arrayrulecolor{black}

\fancypagestyle{plain}{
  \fancyhead{}
  \fancyfoot{}
  \renewcommand{\headrulewidth}{0pt}
}

\usepackage{algorithm}
\usepackage[noend]{algorithmic}

%%% Clear Header %%%%%%%%%%%%%%%%%%%%%%%%%%%%%%%%%%%%%%%%%%%%%%%%%%%%%%%%%%%%%%%%%%
% Clear Header Style on the Last Empty Odd pages
\makeatletter

\def\cleardoublepage{\clearpage\if@twoside \ifodd\c@page\else%
  \hbox{}%
  \thispagestyle{empty}%              % Empty header styles
  \newpage%
  \if@twocolumn\hbox{}\newpage\fi\fi\fi}

\makeatother
 
%%%%%%%%%%%%%%%%%%%%%%%%%%%%%%%%%%%%%%%%%%%%%%%%%%%%%%%%%%%%%%%%%%%%%%%%%%%%%%% 
% Prints your review date and 'Draft Version' (From Josullvn, CS, CMU)
\newcommand{\reviewtimetoday}[2]{\special{!userdict begin
    /bop-hook{gsave 20 710 translate 45 rotate 0.8 setgray
      /Times-Roman findfont 12 scalefont setfont 0 0   moveto (#1) show
      0 -12 moveto (#2) show grestore}def end}}
% You can turn on or off this option.
% \reviewtimetoday{\today}{Draft Version}
%%%%%%%%%%%%%%%%%%%%%%%%%%%%%%%%%%%%%%%%%%%%%%%%%%%%%%%%%%%%%%%%%%%%%%%%%%%%%%% 

\newenvironment{maxime}[1]
{
\vspace*{0cm}
\hfill
\begin{minipage}{0.5\textwidth}%
%\rule[0.5ex]{\textwidth}{0.1mm}\\%
\hrulefill $\:$ {\bf #1}\\
%\vspace*{-0.25cm}
\it 
}%
{%

\hrulefill
\vspace*{0.5cm}%
\end{minipage}
}

\let\minitocORIG\minitoc
\renewcommand{\minitoc}{\minitocORIG \vspace{1.5em}}

\usepackage{multirow}
\usepackage{slashbox}

\newenvironment{bulletList}%
{ \begin{list}%
	{$\bullet$}%
	{\setlength{\labelwidth}{25pt}%
	 \setlength{\leftmargin}{30pt}%
	 \setlength{\itemsep}{\parsep}}}%
{ \end{list} }

\renewcommand{\epsilon}{\varepsilon}

% centered page environment

\newenvironment{vcenterpage}
{\newpage\vspace*{\fill}\fancyhf{}\renewcommand{\headrulewidth}{0pt}}
{\vspace*{\fill}\par\pagebreak}


% Marcador visible de escritura; no eliminar hasta completar su contenido.
\DeclareRobustCommand{\pendiente}[1]{\textcolor{red}{[Pendiente: #1]}}
\pdfstringdefDisableCommands{\def\pendiente#1{Pendiente: #1}}
\newcolumntype{Y}{>{\raggedright\arraybackslash}X}
\makeindex

\begin{document}
% Evitar que la portada sin numeración comparta el destino PDF de la página 1.
\hypersetup{pageanchor=false}
\maketitle
\frontmatter
\hypersetup{pageanchor=true}
% GUÍA: Conservar el texto y formato institucional de ambas cartas, definidos en la clase.
\makeapprovalletter
\makecommitmentletter

\clearpage
\begin{center}{\bfseries Abstract}\end{center}
The visual differentiation of pigmented skin lesions and the identification of melanoma represent a clinical challenge within dermoscopic analysis, due to the morphological similarity between malignant neoplasms and benign lesions. In response to this need, deep learning models based on convolutional neural networks have been established as a computational alternative to assist dermoscopic assessment and guide the prioritization of lesions subject to histopathological confirmation.

This work presents a computational diagnostic support system for the classification of dermoscopic skin lesion images across the seven diagnostic categories of the HAM10000 dataset. The development incorporates a four-stage preprocessing pipeline with noise attenuation, morphological hair removal, and adaptive contrast equalization, together with the comparison of convolutional architectures under a lesion-grouped partition, external validation on the ISIC 2018 dataset, and a probability calibration scheme with configurable decision modes.

The solution was implemented as a web platform that runs the preprocessing and inference core directly in the browser using Web Workers, OpenCV.js, and ONNX Runtime Web, together with a graphical interface developed with React and TypeScript. This approach makes it possible to analyze individual photographs or image batches, inspect the probability distribution across diagnostic categories, and export reports in standardized formats, without sending images to external servers or relying on dedicated computing infrastructure.

The application seeks to facilitate access to computational tools for dermoscopic classification and serve as an analysis and experimentation environment for students, researchers, and healthcare professionals, operating exclusively as a support resource that does not perform definitive diagnoses or replace evaluation by a dermatologist.

\vspace{1em}
\noindent\textbf{Keywords:} melanoma, deep learning, convolutional neural networks, dermoscopy, web application.

\clearpage
\begin{center}{\bfseries Resumen}\end{center}
La diferenciación visual de lesiones cutáneas pigmentadas y la identificación del melanoma representan un reto clínico dentro del análisis dermatoscópico, debido a la similitud morfológica entre neoplasias malignas y lesiones benignas. Frente a esta necesidad, los modelos de aprendizaje profundo basados en redes neuronales convolucionales se han consolidado como una alternativa computacional para asistir la valoración dermatoscópica y orientar la priorización de lesiones sujetas a confirmación histopatológica.

Este trabajo presenta un sistema computacional de apoyo al diagnóstico para la clasificación de imágenes dermatoscópicas de lesiones cutáneas en las siete categorías diagnósticas del conjunto HAM10000. El desarrollo incorpora una cadena de preprocesamiento de cuatro etapas con atenuación de ruido, remoción morfológica de vello y ecualización adaptativa de contraste, junto con la comparación de arquitecturas convolucionales bajo una partición agrupada por lesión, la validación externa en el conjunto ISIC 2018 y un esquema de calibración de probabilidades con modos de decisión configurables.

La solución fue implementada como una plataforma web que ejecuta el núcleo de preprocesamiento e inferencia directamente en el navegador mediante Web Workers, OpenCV.js y ONNX Runtime Web, junto con una interfaz gráfica desarrollada con React y TypeScript. Este enfoque permite analizar fotografías individuales o lotes de imágenes, inspeccionar la distribución de probabilidades por categoría diagnóstica y exportar reportes en formatos estandarizados, sin enviar las imágenes a servidores externos ni depender de infraestructura de cómputo dedicada.

El aplicativo busca facilitar el acceso a herramientas computacionales de clasificación dermatoscópica y servir como un entorno de análisis y experimentación para estudiantes, investigadores y profesionales del área, operando exclusivamente como un recurso de apoyo que no realiza diagnósticos definitivos ni sustituye la evaluación del dermatólogo.

\vspace{1em}
\noindent\textbf{Palabras clave:} melanoma, aprendizaje profundo, redes neuronales convolucionales, dermatoscopia, aplicativo web.

\tableofcontents
\listoftables
\listoffigures

\mainmatter
\chapter{Introducción}

El melanoma cutáneo es la variante de cáncer de piel con mayor capacidad de diseminación metastásica, y su pronóstico clínico depende de su identificación en fases iniciales~\cite{amaral2025esmo}. Según las estimaciones mundiales de GLOBOCAN para 2024 (publicadas en julio de 2026), esta neoplasia registró 337.969 casos nuevos y 58.398 muertes en el mundo~\cite{sung2026globocan,iarc2026melanoma}. En Colombia, la Cuenta de Alto Costo informó que el 67,4\,\% de los casos diagnosticados en el periodo 2024 se identificó en estadios tempranos I y II, dentro de un reporte preliminar de 11.064 casos atendidos~\cite{cac2025melanoma}. Aunque la dermatoscopia realizada por especialistas mejora la precisión diagnóstica antes de la confirmación histopatológica obligatoria~\cite{amaral2025esmo}, la disponibilidad de herramientas tecnológicas de apoyo sigue siendo limitada fuera de entornos clínicos especializados~\cite{un2024sdg}.

El aprendizaje profundo ha cobrado relevancia en el análisis y clasificación de imágenes médicas~\cite{goodfellow2016deep,esteva2017dermatologist}. Las redes neuronales convolucionales (CNN) permiten extraer características morfológicas de imágenes dermatoscópicas y diferenciar entre lesiones benignas y malignas como apoyo a la valoración de los especialistas~\cite{esteva2017dermatologist,sariates2025}. Este tipo de clasificadores se ha evaluado incluso en estudios clínicos prospectivos como PROVE-AI (603 lesiones)~\cite{marchetti2023proveai} y en cohortes multicéntricas de 1.910 lesiones en ocho hospitales universitarios~\cite{heinlein2024multicenter}. De forma complementaria, el aprendizaje por transferencia (\textit{transfer learning}) adapta modelos preentrenados en ImageNet~\cite{deng2009imagenet} al dominio de las lesiones cutáneas, reduciendo los datos y el tiempo de entrenamiento requeridos~\cite{zhang2024dive,cherukuri2025,bello2024densenet}.

No obstante, persiste una brecha ingenieril y metodológica entre los modelos experimentales y su implementación en soluciones funcionales para el usuario~\cite{farias2026,almamun2025}. En el plano metodológico, el conjunto de referencia HAM10000 reúne 10.015 imágenes de 7.470 lesiones (con 1.956 lesiones fotografiadas más de una vez)~\cite{tschandl2018ham10000}, por lo que dividir los datos por imagen reparte fotografías de una misma lesión entre entrenamiento y prueba e infla las métricas por fuga de información~\cite{kapoor2023leakage,abhishek2025quality}. En el plano de despliegue, gran parte de las propuestas se limita a entornos de laboratorio y presenta barreras de usabilidad, privacidad o dependencia de servidores dedicados~\cite{farias2026,ning2024onnxweb}. El problema exige tanto entrenar bajo una partición libre de fuga como conectar el clasificador con una interfaz web desplegable, verificada y accesible~\cite{almamun2025,wang2025browser}.

El presente trabajo desarrolla un sistema de apoyo al diagnóstico de melanoma sobre imágenes dermatoscópicas que integra CNN y \textit{transfer learning} con un aplicativo web funcional desplegado en \pendiente{URL}. El sistema cubre el preprocesamiento de imágenes en cuatro etapas, el entrenamiento y evaluación de ocho configuraciones en 224 y 384\,px bajo una partición agrupada por lesión sobre HAM10000~\cite{tschandl2018ham10000}, la selección preregistrada de MobileNetV3Large~\cite{howard2019mobilenetv3} y su validación externa en el conjunto International Skin Imaging Collaboration (ISIC) 2018~\cite{codella2019isic2018}. Adicionalmente, incorpora la calibración de probabilidades~\cite{guo2017calibration} con modos de decisión configurables y su despliegue en una plataforma web que ejecuta el análisis en el navegador mediante OpenCV.js y ONNX Runtime Web~\cite{ning2024onnxweb}, donde los usuarios pueden cargar imágenes dermatoscópicas y obtener una predicción preliminar de apoyo sin enviarlas a servidores externos. El aporte del proyecto articula un protocolo de clasificación libre de fuga por lesión con una herramienta web accesible, conectando la experimentación en aprendizaje profundo con su despliegue práctico como apoyo a la valoración dermatoscópica~\cite{amaral2025esmo,farias2026}.

El resto del documento se estructura en ocho capítulos adicionales. El Capítulo~2 presenta la descripción del problema, la pregunta de investigación, la sistematización, los objetivos, la justificación vinculada al Objetivo de Desarrollo Sostenible (ODS) 3~\cite{un2024sdg} y las delimitaciones. El Capítulo~3 desarrolla el marco teórico sobre lesiones pigmentadas, dermatoscopia, CNN, \textit{transfer learning}, desbalance y métricas~\cite{maierhein2024metrics}, junto con los trabajos relacionados. El Capítulo~4 describe el conjunto HAM10000~\cite{tschandl2018ham10000}, la partición agrupada por lesión, el preprocesamiento, el aumento en línea y los pesos de clase; a continuación, el Capítulo~5 detalla las arquitecturas, la receta en dos fases, el experimento de fuga por unidad de partición, la regla de selección y la calibración~\cite{guo2017calibration}. Por su parte, el Capítulo~6 documenta la arquitectura del aplicativo web~\cite{ning2024onnxweb}, los requisitos, la interfaz, el manual de usuario y la verificación técnica~\cite{wang2025browser}. Finalmente, el Capítulo~7 expone los resultados en HAM10000 e ISIC 2018~\cite{codella2019isic2018}, la discusión y las limitaciones; el Capítulo~8 sintetiza las conclusiones por objetivo específico y el Capítulo~9 plantea el trabajo futuro~\cite{vasey2022decideai,daneshjou2022disparities,collins2024tripodai}.

\chapter{Descripción del problema}

Este capítulo plantea el problema de investigación, su formulación y las cuatro preguntas de sistematización que guían el desarrollo experimental. También expone el objetivo general y los cuatro objetivos específicos aprobados en el anteproyecto para estructurar las etapas del sistema. Finalmente, presenta la justificación sanitaria, técnica y de ingeniería junto con las delimitaciones del prototipo en nivel de madurez tecnológica (TRL) 4~\cite{amaral2025esmo,vasey2022decideai}.

\section{Formulación}

El melanoma cutáneo registró 337.969 casos nuevos y 58.398 muertes estimadas en el mundo durante 2024~\cite{sung2026globocan,iarc2026melanoma}. En la práctica clínica, la dermatoscopia mejora la precisión diagnóstica de los especialistas y el diagnóstico definitivo exige confirmación histopatológica~\cite{amaral2025esmo}. Aunque se han entrenado clasificadores sobre las 10.015 imágenes de HAM10000~\cite{tschandl2018ham10000}, su evaluación se distorsiona cuando imágenes de una misma lesión quedan en entrenamiento y prueba~\cite{kapoor2023leakage,abhishek2025quality}. Además, gran parte de los desarrollos permanecen como experimentos de laboratorio sin integrarse en aplicativos web accesibles~\cite{farias2026}. A partir de estas necesidades metodológicas y de despliegue, el anteproyecto definió la siguiente pregunta de investigación:

\begin{quote}
¿Cómo desarrollar un sistema de apoyo al diagnóstico de melanoma, basado en redes neuronales convolucionales y técnicas de \textit{transfer learning}, integrado en un aplicativo web que permita la clasificación de imágenes de lesiones cutáneas y obtenga una predicción de riesgo de forma accesible?
\end{quote}

\section{Sistematización}

Para descomponer la pregunta de investigación, el problema se estructuró en cuatro subpreguntas alineadas con las etapas del proyecto. Estas subpreguntas abarcan la gestión de las 10.015 imágenes de HAM10000~\cite{tschandl2018ham10000}, la implementación de arquitecturas con aprendizaje por transferencia~\cite{zhang2024dive}, la selección del clasificador~\cite{maierhein2024metrics} y su despliegue web~\cite{ning2024onnxweb}. A continuación se transcriben textualmente las cuatro subpreguntas formuladas en el anteproyecto:

\begin{enumerate}
  \item ¿Cómo gestionar una base de datos compuesta por imágenes cutáneas que sea etiquetada y que permita la identificación de melanoma?
  \item ¿Cuáles arquitecturas de aprendizaje profundo se pueden implementar para la clasificación de melanoma a partir de imágenes cutáneas?
  \item ¿Cómo seleccionar el mejor modelo de aprendizaje profundo para la identificación de me\-la\-no\-ma?
  \item ¿Cómo integrar el modelo de clasificación seleccionado en un aplicativo web funcional que permita a usuarios cargar imágenes cutáneas y obtener un resultado de identificación?
\end{enumerate}

\section{Objetivos}

Los objetivos del trabajo establecen el propósito general y las cuatro metas operativas correspondientes a las preguntas de sistematización de la Sección~2.2. Su cumplimiento abarca desde el procesamiento del conjunto HAM10000~\cite{tschandl2018ham10000} hasta la verificación del aplicativo web con inferencia local~\cite{ning2024onnxweb,wang2025browser}. A continuación se transcriben de forma textual el objetivo general y los cuatro objetivos específicos del anteproyecto aprobado.

\subsection{Objetivo general}

Desarrollar un sistema de detección de cáncer de piel mediante técnicas de \textit{deep learning}, integrado en un aplicativo web, analizando imágenes de lesiones cutáneas y generando una predicción como apoyo al diagnóstico preventivo.

\subsection{Objetivos específicos}

\begin{enumerate}
  \item Gestionar una base de datos compuesta por imágenes cutáneas etiquetadas aplicando limpieza, preprocesamiento y aumento de datos para el entrenamiento de métodos de aprendizaje profundo.
  \item Implementar arquitecturas basadas en redes neuronales convolucionales (CNN) para la clasificación de melanoma a partir de imágenes cutáneas.
  \item Seleccionar el mejor modelo de aprendizaje profundo para la identificación de melanoma mediante la evaluación de los métodos con métricas estándar de clasificación.
  \item Integrar el modelo seleccionado en un aplicativo web funcional que permita a los usuarios cargar imágenes cutáneas y obtener un resultado de identificación de melanoma.
\end{enumerate}

\section{Justificación}

Desde la perspectiva sanitaria, las 58.398 muertes y los 337.969 casos nuevos estimados en el mundo en 2024 evidencian la carga del melanoma cutáneo~\cite{sung2026globocan,iarc2026melanoma}. En Colombia, la Cuenta de Alto Costo informó que el 67,4\,\% de los casos diagnosticados en el periodo 2024 correspondió a estadios I y II~\cite{cac2025melanoma}. En ese mismo reporte nacional se registraron 11.064 casos atendidos de carácter preliminar, previos a auditoría~\cite{cac2025melanoma}. La finalidad de este prototipo se alinea conceptualmente con el ODS 3 (Salud y bienestar), que plantea fortalecer la prevención y la capacidad de los sistemas sanitarios~\cite{un2024sdg}. Esta relación expresa la orientación del sistema como herramienta de apoyo y no representa una reducción de mortalidad ni un impacto clínico demostrado~\cite{vasey2022decideai}.

En el plano técnico y metodológico, la literatura ha documentado que la fuga de información produce estimaciones optimistas en 17 campos científicos~\cite{kapoor2023leakage}. En dermatología computacional, el conjunto HAM10000 incluye 1.956 lesiones con múltiples fotografías, las cuales suman 4.501 de sus 10.015 imágenes~\cite{tschandl2018ham10000}. Auditorías recientes sobre derivados de HAM10000 demostraron que partir los datos por imagen reparte fotografías de una misma lesión entre entrenamiento y prueba~\cite{abhishek2025quality}. Por ello, este trabajo aplicó una partición por lesión (5.336 lesiones en entrenamiento, 1.067 en validación y 1.067 en prueba) y midió la diferencia frente a partir por imagen. Además, el modelo seleccionado se evaluó sin reajustar sus umbrales sobre 1.512 imágenes de ISIC 2018 sin lesiones compartidas con el entrenamiento~\cite{codella2019isic2018}.

Desde la ingeniería de software, la integración del clasificador mediante ONNX Runtime Web ejecuta el preprocesamiento y la inferencia directamente en el navegador del cliente~\cite{ning2024onnxweb}. Esta arquitectura evita mantener servidores con unidad de procesamiento gráfico (GPU) y procesa cada imagen en 59\,ms en promedio en Chromium sobre unidad central de procesamiento (CPU). De igual forma, las imágenes dermatoscópicas no abandonan el equipo del usuario durante el análisis~\cite{ning2024onnxweb}. Frente al planteamiento inicial del anteproyecto, se precisa que el sistema no está destinado a que la población general evalúe lesiones sin intervención clínica. El aplicativo requiere imágenes adquiridas con dermatoscopio y entrega una predicción de apoyo que no reemplaza el criterio médico ni la confirmación histopatológica~\cite{amaral2025esmo}.

\section{Delimitaciones y alcance}

El sistema desarrollado corresponde a un prototipo validado en entorno de laboratorio (TRL 4) mediante evaluación retrospectiva sobre conjuntos públicos~\cite{tschandl2018ham10000,codella2019isic2018}. Aunque el anteproyecto preveía una ejecución restringida a un entorno local, el aplicativo final se desplegó como sitio web estático en Vercel (\pendiente{URL}). Este cambio de alcance fue posible porque el modelo exportado a ONNX (12,9\,MB) y el preprocesamiento en OpenCV.js se ejecutan íntegramente en el navegador del usuario~\cite{ning2024onnxweb}. De este modo, el despliegue público no requiere infraestructura de cómputo en servidor ni transmisión de imágenes por la red.

El alcance del modelo se circunscribe a imágenes dermatoscópicas clasificadas en las siete categorías pigmentadas definidas en HAM10000~\cite{tschandl2018ham10000}. El clasificador no fue entrenado ni evaluado con fotografías clínicas tomadas con teléfonos móviles ni incluye metadatos del paciente como edad, sexo o localización anatómica. El proyecto no comprende ensayos clínicos prospectivos con pacientes~\cite{marchetti2023proveai,heinlein2024multicenter}, estudios de integración en flujos asistenciales según la guía DECIDE-AI~\cite{vasey2022decideai} ni módulos de interoperabilidad hospitalaria. En consecuencia, toda salida del aplicativo debe interpretarse como información de apoyo sujeta a valoración dermatológica e histopatológica~\cite{amaral2025esmo}.

\chapter{Componente teórico}

Este capítulo expone los fundamentos clínicos y computacionales que sustentan el sistema de apoyo al diagnóstico sobre imágenes dermatoscópicas~\cite{amaral2025esmo,zhang2024dive}. En su primera sección desarrolla las categorías de lesiones cutáneas de HAM10000~\cite{tschandl2018ham10000}, los principios de las CNN, la transferencia de aprendizaje, el desbalance y las métricas de validación~\cite{maierhein2024metrics}. En la segunda sección analiza los trabajos relacionados a partir de sus protocolos de partición, sus métricas y su alcance clínico~\cite{kapoor2023leakage,barata2023reinforcement}.

\section{Marco teórico}

El marco teórico articula los conceptos médicos y algorítmicos requeridos para clasificar imá\-ge\-nes dermatoscópicas en siete categorías diagnósticas~\cite{tschandl2018ham10000}. Se describen primero las lesiones pigmentadas y la técnica de adquisición dermatoscópica utilizada en la práctica clínica~\cite{amaral2025esmo}. Posteriormente se formalizan las operaciones de las CNN, el ajuste fino desde pesos preentrenados, el tratamiento del desbalance y las ecuaciones de evaluación y calibración~\cite{zhang2024dive,maierhein2024metrics,guo2017calibration}.

\subsection{Lesiones cutáneas y melanoma}

El conjunto de datos HAM10000 agrupa las lesiones pigmentadas más frecuentes en la práctica dermatoscópica dentro de siete clases diagnósticas~\cite{tschandl2018ham10000}. La Tabla~\ref{tab:clases} presenta el código identificador, la denominación en español, el comportamiento biológico y una descripción resumida de cada una de estas siete categorías~\cite{tschandl2018ham10000}. En ella se observa la coexistencia de lesiones benignas, entidades premalignas o malignas \textit{in situ} y neoplasias malignas invasivas como el melanoma (\texttt{mel}) y el carcinoma basocelular (\texttt{bcc})~\cite{tschandl2018ham10000}.

\begin{table}[H]
  \centering
  \caption{Clases diagnósticas del conjunto de datos HAM10000~\cite{tschandl2018ham10000}.}
  \label{tab:clases}
  \begin{tabularx}{\textwidth}{l >{\hsize=0.9\hsize\linewidth=\hsize}Y >{\hsize=0.65\hsize\linewidth=\hsize}Y >{\hsize=1.45\hsize\linewidth=\hsize}Y}
    \toprule
    Código & Nombre & Naturaleza & Descripción \\
    \midrule
    \texttt{akiec} & Queratosis actínica y enfermedad de Bowen & Premaligna / maligna \textit{in situ} & Variantes no invasivas de carcinoma escamocelular asociadas a daño actínico crónico~\cite{tschandl2018ham10000}. \\
    \texttt{bcc} & Carcinoma basocelular & Maligna & Neoplasia epitelial maligna de crecimiento lento que rara vez produce metástasis~\cite{tschandl2018ham10000}. \\
    \texttt{bkl} & Lesiones benignas tipo queratosis & Benigna & Agrupa queratosis seborreica, lentigo solar y queratosis tipo liquen plano~\cite{tschandl2018ham10000}. \\
    \texttt{df} & Dermatofibroma & Benigna & Proliferación dérmica benigna reactiva o fibrocítica de aspecto nodular~\cite{tschandl2018ham10000}. \\
    \texttt{mel} & Melanoma & Maligna & Neoplasia maligna originada en los melanocitos con capacidad metastásica~\cite{tschandl2018ham10000,amaral2025esmo}. \\
    \texttt{nv} & Nevo melanocítico & Benigna & Proliferación benigna de melanocitos que presenta múltiples variantes morfológicas~\cite{tschandl2018ham10000}. \\
    \texttt{vasc} & Lesiones vasculares & Benigna & Incluye angiomas, angioqueratomas y granulomas piógenos de coloración rojiza o violácea~\cite{tschandl2018ham10000}. \\
    \bottomrule
  \end{tabularx}
\end{table}

El melanoma (\texttt{mel}) es una neoplasia maligna derivada de los melanocitos que ocasionó 58.398 muertes estimadas en el mundo en 2024~\cite{iarc2026melanoma,tschandl2018ham10000}. Su identificación en estadios tempranos I y II, que en Colombia representaron el 67,4\,\% de los diagnósticos del periodo 2024~\cite{cac2025melanoma}, condiciona el manejo quirúrgico y el pronóstico~\cite{amaral2025esmo}. En dermatoscopia, el melanoma comparte rasgos visuales de pigmentación y asimetría con los nevos melanocíticos (\texttt{nv}) y con las lesiones benignas tipo queratosis (\texttt{bkl})~\cite{tschandl2018ham10000}.

Dentro de las categorías restantes, la clase \texttt{bkl} reúne tres entidades benignas distintas: la queratosis seborreica, el lentigo solar y la queratosis tipo liquen plano~\cite{tschandl2018ham10000}. Estas lesiones se agrupan en una sola etiqueta porque muestran patrones dermatoscópicos similares e inflamación regresiva~\cite{tschandl2018ham10000}. Por su parte, la clase \texttt{akiec} agrupa la queratosis actínica y el carcinoma intraepitelial o enfermedad de Bowen, que constituyen precursores o formas superficiales no invasivas de carcinoma escamocelular~\cite{tschandl2018ham10000}. Las tres categorías restantes corresponden al carcinoma basocelular (\texttt{bcc}), el dermatofibroma (\texttt{df}) y las lesiones vasculares (\texttt{vasc})~\cite{tschandl2018ham10000}.

\subsection{Dermatoscopia}

La dermatoscopia es una técnica de microscopía de epiluminiscencia no invasiva que emplea iluminación magnificada con luz polarizada o líquido de inmersión~\cite{amaral2025esmo,tschandl2018ham10000}. Este procedimiento reduce la reflexión superficial de la epidermis y permite visualizar estructuras de pigmento y patrones vasculares en la epidermis profunda y la dermis papilar~\cite{tschandl2018ham10000}. Seg\'{u}n la gu\'{\i}a cl\'{\i}nica de la Sociedad Europea de Oncolog\'{\i}a M\'{e}dica (ESMO), la dermatoscopia realizada por un m\'{e}dico con experiencia mejora la precisi\'{o}n diagn\'{o}stica~\cite{amaral2025esmo}.

Pese a su utilidad clínica, el examen dermatoscópico no establece por sí solo el diagnóstico definitivo de malignidad~\cite{amaral2025esmo}. Ante una lesión sospechosa, la confirmación de melanoma requiere la escisión de la lesión y su posterior estudio histopatológico~\cite{amaral2025esmo}. Adicionalmente, las imágenes dermatoscópicas suelen presentar artefactos de adquisición como burbujas, variaciones de iluminación y oclusión por vello cutáneo~\cite{lee1997dullrazor,wang2024hair}. Estudios sobre HAM10000 indican que la presencia abundante de vello reduce el rendimiento de los clasificadores automáticos, con mayor impacto en las lesiones vasculares (\texttt{vasc})~\cite{wang2024hair}.

\subsection{Redes neuronales convolucionales}

Las CNN procesan imágenes en forma de tensores tridimensionales de altura, anchura y canales de color mediante capas sucesivas con pesos compartidos~\cite{goodfellow2016deep,zhang2024dive}. En cada capa convolucional, un conjunto de filtros bidimensionales recorre la entrada para extraer mapas de características locales como bordes, texturas y patrones de color~\cite{goodfellow2016deep}. Dado un mapa de entrada bidimensional $I$ y un núcleo o filtro $K$ de tamaño $M \times N$, la operación de convolución discreta en la posición $(i, j)$ se expresa en la Ecuación~\ref{eq:convolucion}~\cite{goodfellow2016deep}:

\begin{equation}
  S(i, j) = (I * K)(i, j) = \sum_{m=1}^{M} \sum_{n=1}^{N} I(i + m, j + n)\,K(m, n)
  \label{eq:convolucion}
\end{equation}

Sobre los mapas resultantes se aplica una función de activación no lineal, como la unidad lineal rectificada (ReLU), definida por $f(x) = \max(0, x)$~\cite{goodfellow2016deep,zhang2024dive}. A continuación, las capas de agrupamiento (\textit{pooling}) reducen la resolución espacial de los mapas de características para disminuir el número de parámetros y otorgar invarianza a traslaciones pequeñas~\cite{goodfellow2016deep}. Al final de la red, el agrupamiento promedio global resume cada canal espacial en un escalar antes de las capas densas y la regularización con \textit{dropout}~\cite{zhang2024dive}. Para un problema de $C = 7$ clases, la capa de salida aplica la función \textit{softmax} sobre el vector de \textit{logits} $z = (z_1, \dots, z_C)$ según la Ecuación~\ref{eq:softmax}~\cite{goodfellow2016deep}:

\begin{equation}
  p_c = \frac{\exp(z_c)}{\sum_{k=1}^{C} \exp(z_k)}, \quad c \in \{1, \dots, C\}
  \label{eq:softmax}
\end{equation}

En este trabajo se evaluaron cinco familias de arquitecturas convolucionales preentrenadas junto con variantes recurrentes. ResNet50 introduce conexiones residuales de atajo que suman la entrada de un bloque a su salida para facilitar la propagación del gradiente en redes de 50 capas~\cite{he2016resnet}. DenseNet121 conecta cada capa con las siguientes dentro de un bloque denso mediante concatenación de características, favoreciendo la reutilización de representaciones con 7,30 millones de parámetros~\cite{huang2017densenet}. MobileNetV3Large combina convoluciones separables en profundidad, bloques residuales invertidos con atención de compresión y excitación y búsqueda de arquitectura neuronal, con un total de 3,24 millones de parámetros~\cite{howard2019mobilenetv3}.

Las dos arquitecturas preentrenadas restantes incorporan rediseños posteriores de las convoluciones modernas. EfficientNetV2-S combina bloques convolucionales fusionados en las primeras etapas con un escalado compuesto de profundidad, anchura y resolución, sumando 20,66 millones de parámetros~\cite{tan2021efficientnetv2}. ConvNeXt-Tiny moderniza el diseño de una red residual estándar mediante convoluciones en profundidad de $7 \times 7$ y bloques invertidos con 28,02 millones de parámetros~\cite{liu2022convnext}. Por último, añadir una capa de memoria a corto y largo plazo bidireccional (BiLSTM) sobre los mapas de características modela dependencias secuenciales en ambas direcciones~\cite{zhang2024dive}.

\subsection{Aprendizaje por transferencia y ajuste fino}

El aprendizaje por transferencia reutiliza los parámetros de una red preentrenada sobre un conjunto de gran escala, como ImageNet~\cite{deng2009imagenet}, para inicializar un clasificador en un dominio con menos muestras~\cite{zhang2024dive}. Las primeras capas convolucionales aprenden detectores generales de bordes, gradientes cromáticos y texturas que resultan transferibles a imágenes dermatoscópicas~\cite{zhang2024dive}. Para adaptar la red a HAM10000, se reemplaza el clasificador original de 1.000 categorías de ImageNet~\cite{deng2009imagenet} por una cabeza con siete salidas según la Tabla~\ref{tab:clases}~\cite{zhang2024dive}.

El procedimiento de ajuste fino (\textit{fine-tuning}) se ejecuta habitualmente en dos etapas sucesivas para preservar las representaciones preentrenadas~\cite{zhang2024dive}. En la primera fase se congelan los pesos de la base convolucional y se entrena únicamente la cabeza de clasificación inicializada al azar~\cite{zhang2024dive}. En la segunda fase se descongelan las capas convolucionales con una tasa de aprendizaje reducida para ajustar los filtros a las estructuras específicas de las lesiones cutáneas~\cite{zhang2024dive}. Como tendencia reciente, se han creado modelos fundacionales dermatológicos como PanDerm, preentrenado con más de dos millones de imágenes de once instituciones y cuatro modalidades~\cite{yan2025panderm}.

\subsection{Desbalance de clases}

Un conjunto de datos presenta desbalance de clases cuando la frecuencia de las categorías difiere en uno o más órdenes de magnitud~\cite{hu2024multiscale}. En HAM10000, los nevos melanocíticos (\texttt{nv}) concentran 6.705 de las 10.015 imágenes (67,0\,\%), mientras que el dermatofibroma (\texttt{df}) cuenta con 115 imágenes y las lesiones vasculares (\texttt{vasc}) con 142~\cite{tschandl2018ham10000}. Si una red se entrena minimizando la entropía cruzada estándar sin corrección, el gradiente queda dominado por la clase mayoritaria y el clasificador tiende a subestimar las clases minoritarias~\cite{hu2024multiscale}.

En la literatura sobre clasificación dermatoscópica se han comparado tres familias principales de estrategias para mitigar este problema~\cite{hu2024multiscale}. La primera consiste en el remuestreo de datos, ya sea mediante sobremuestreo de las categorías minoritarias o submuestreo de la clase mayoritaria~\cite{hu2024multiscale}. La segunda asigna pesos por clase $w_c$ en la función de pérdida de entropía cruzada para penalizar en mayor proporción los errores sobre las clases menos frecuentes~\cite{hu2024multiscale}. La tercera combina estas técnicas con regularización de las probabilidades objetivo mediante suavizado de etiquetas (\textit{label smoothing}) para reducir el exceso de confianza del modelo~\cite{hu2024multiscale}.

\subsection{Métricas de evaluación}

En un problema de clasificación con desbalance severo, la proporción global de aciertos (accuracy) resulta insuficiente porque puede alcanzar valores altos acertando únicamente la clase mayoritaria~\cite{maierhein2024metrics}. Por ello, las guías metodológicas recomiendan seleccionar métricas acordes con la distribución de clases y reportar el rendimiento desagregado por categoría~\cite{maierhein2024metrics}. Sean $TP_c$, $FP_c$, $TN_c$ y $FN_c$ los verdaderos positivos, falsos positivos, verdaderos negativos y falsos negativos de la clase $c \in \{1, \dots, C\}$ sobre un total de $N$ muestras. La accuracy global se define en la Ecuación~\ref{eq:accuracy}~\cite{maierhein2024metrics}:

\begin{equation}
  \text{Accuracy} = \frac{1}{N} \sum_{c=1}^{C} TP_c
  \label{eq:accuracy}
\end{equation}

Para una clase específica $c$ como el melanoma (\texttt{mel}), la sensibilidad y la especificidad miden las tasas de verdaderos positivos y verdaderos negativos~\cite{maierhein2024metrics}. Ambas métricas se definen en las Ecuaciones~\ref{eq:sensibilidad} y~\ref{eq:especificidad}~\cite{maierhein2024metrics}:

\begin{equation}
  \text{Sensibilidad}_c = \frac{TP_c}{TP_c + FN_c}
  \label{eq:sensibilidad}
\end{equation}

\begin{equation}
  \text{Especificidad}_c = \frac{TN_c}{TN_c + FP_c}
  \label{eq:especificidad}
\end{equation}

De forma complementaria, la precisión (\textit{precision}) cuantifica la proporción de predicciones positivas de la clase $c$ que resultan verdaderas~\cite{maierhein2024metrics}. Su media armónica con la sensibilidad define el puntaje $\text{F1}_c$, según las Ecuaciones~\ref{eq:precision} y~\ref{eq:f1}~\cite{maierhein2024metrics}:

\begin{equation}
  \text{Precisión}_c = \frac{TP_c}{TP_c + FP_c}
  \label{eq:precision}
\end{equation}

\begin{equation}
  \text{F1}_c = \frac{2 \cdot \text{Precisión}_c \cdot \text{Sensibilidad}_c}{\text{Precisión}_c + \text{Sensibilidad}_c} = \frac{2\,TP_c}{2\,TP_c + FP_c + FN_c}
  \label{eq:f1}
\end{equation}

Para resumir el desempeño multiclase otorgando el mismo peso a las siete categorías con independencia de su frecuencia, se utilizan el Macro-F1 y la balanced accuracy~\cite{maierhein2024metrics}. Esta última corresponde al promedio aritmético de las sensibilidades por clase y fue adoptada como métrica oficial del reto ISIC 2018 Task 3~\cite{codella2019isic2018}. Sus expresiones se presentan en las Ecuaciones~\ref{eq:macrof1} y~\ref{eq:balacc}~\cite{maierhein2024metrics,codella2019isic2018}:

\begin{equation}
  \text{Macro-F1} = \frac{1}{C} \sum_{c=1}^{C} \text{F1}_c
  \label{eq:macrof1}
\end{equation}

\begin{equation}
  \text{Balanced accuracy} = \frac{1}{C} \sum_{c=1}^{C} \text{Sensibilidad}_c
  \label{eq:balacc}
\end{equation}

Además de las métricas dependientes de un punto de corte, la capacidad de discriminación para melanoma se evalúa mediante la curva característica operativa del receptor (ROC)~\cite{maierhein2024metrics}. Esta curva grafica la sensibilidad frente a $1 - \text{Especificidad}_{\texttt{mel}}$ al variar el umbral de decisión $t \in [0, 1]$, y el área bajo la curva (AUC) define el AUC-ROC en la Ecuación~\ref{eq:auc}~\cite{maierhein2024metrics}:

\begin{equation}
  \text{AUC}_{\texttt{mel}} = \int_{0}^{1} \text{Sensibilidad}_{\texttt{mel}}\!\left(\text{FPR}^{-1}(u)\right) du
  \label{eq:auc}
\end{equation}

En el análisis de imágenes médicas es necesario distinguir entre la discriminación del modelo y la calibración de sus probabilidades~\cite{maierhein2024metrics,sambyal2023calibration}. Una red neuronal puede alcanzar una accuracy o un AUC altos y, al mismo tiempo, producir probabilidades sobreconfiadas que no reflejan la frecuencia empírica de acierto~\cite{guo2017calibration,sambyal2023calibration}. El método de \textit{temperature scaling} corrige esta descalibración dividiendo los \textit{logits} $z$ por un escalar $T > 0$ antes de la función \textit{softmax} (Ecuación~\ref{eq:tempscaling})~\cite{guo2017calibration}. El parámetro $T$ se ajusta sobre un conjunto de validación reservado y su comportamiento depende de la confiabilidad de las etiquetas clínicas~\cite{penso2024calibration}.

\begin{equation}
  p_{\text{cal}, c} = \frac{\exp(z_c / T)}{\sum_{k=1}^{C} \exp(z_k / T)}
  \label{eq:tempscaling}
\end{equation}

\section{Trabajos relacionados}

En la clasificación multiclase sobre HAM10000~\cite{tschandl2018ham10000}, diversos estudios han explorado modificaciones arquitectónicas y esquemas de balanceo. Hu et al.~\cite{hu2024multiscale} compararon ponderación de clases, \textit{label smoothing} y remuestreo en HAM10000 sin vello, reportando 94,0\,\% de accuracy y 99,3\,\% de AUC (y 89,8\,\% en ISIC 2019). De igual modo, se han evaluado ensambles de CNN y transformadores~\cite{sariates2025}, marcos como EnsembleSkinNet~\cite{cherukuri2025}, ajustes de DenseNet121~\cite{bello2024densenet} y redes de bajo costo~\cite{almamun2025}. Sin embargo, en varios de estos reportes la partición independiente por lesión no está verificada y el uso de accuracy global no es equiparable al Macro-F1 en presencia de desbalance~\cite{maierhein2024metrics,abhishek2025quality}.

Una segunda línea de investigación aborda la asimetría en el costo clínico de los errores diagnósticos al priorizar la detección de melanoma~\cite{barata2023reinforcement}. Barata et al.~\cite{barata2023reinforcement} entrenaron un sistema con HAM10000 y lo evaluaron sobre 1.511 imágenes del conjunto de prueba de ISIC 2018 incorporando recompensas diferenciadas mediante aprendizaje por refuerzo. Al penalizar los falsos negativos en lesiones malignas, su modelo incrementó la sensibilidad para melanoma de 61,4\,\% a 79,5\,\%~\cite{barata2023reinforcement}. Este estudio constituye el antecedente más cercano a los umbrales de este trabajo, aunque usa otro método de decisión y evalúa 1.511 imágenes frente a 1.512~\cite{barata2023reinforcement,codella2019isic2018}.

En el ámbito de la validación clínica con pacientes, los estudios prospectivos han examinado el comportamiento de clasificadores dermatoscópicos fuera de los bancos retrospectivos. En el estudio prospectivo PROVE-AI, Marchetti et al.~\cite{marchetti2023proveai} evaluaron 603 lesiones seleccionadas para biopsia en 435 participantes (96\,\% de piel blanca), incluidas 95 lesiones de melanoma. Con un umbral fijado para 95\,\% de sensibilidad, el modelo obtuvo 96,8\,\% de sensibilidad y 37,4\,\% de especificidad, requiriendo ensayos aleatorizados antes de su adopción~\cite{marchetti2023proveai}. En otro estudio sobre 1.910 lesiones de 1.716 pacientes en ocho hospitales alemanes, Heinlein et al.~\cite{heinlein2024multicenter} reportaron 92,1\,\% de sensibilidad y 67,3\,\% de especificidad (73,4\,\% y 82,8\,\% en dermatólogos).

En cuanto al escalamiento de representaciones y la diversidad poblacional, la literatura reciente ha evidenciado tanto el avance de los modelos fundacionales como las disparidades de desempeño. Yan et al.~\cite{yan2025panderm} desarrollaron PanDerm mediante preentrenamiento multimodal sobre más de dos millones de imágenes procedentes de once instituciones y cuatro modalidades, evaluándolo en 28 tareas clínicas. Por su parte, Daneshjou et al.~\cite{daneshjou2022disparities} demostraron con fotografías clínicas que los modelos rinden peor en pieles oscuras y que el ajuste con datos diversos reduce esa brecha. Ambos trabajos ponen de relieve las limitaciones de entrenar clasificadores únicamente sobre colecciones retrospectivas de poblaciones de piel clara~\cite{daneshjou2022disparities,yan2025panderm}.

Desde el punto de vista de la calidad de los datos y el despliegue en software, las auditorías metodológicas han identificado riesgos críticos de contaminación entre particiones~\cite{kapoor2023leakage,abhishek2025quality}. Kapoor y Narayanan~\cite{kapoor2023leakage} sistematizaron una taxonomía de ocho tipos de fuga de información y documentaron su presencia en 17 campos científicos como causa de sobreestimación del rendimiento. En dermatología, Abhishek et al.~\cite{abhishek2025quality} comprobaron que DermaMNIST reparte imágenes de una misma lesión entre particiones y hallaron 18 pares duplicados en HAM10000 con distinto \texttt{lesion\_id}. En paralelo, los aplicativos web dermatológicos~\cite{farias2026} que ejecutan modelos en el navegador deben verificar paridad numérica, latencia y consumo de memoria~\cite{ning2024onnxweb,wang2025browser}. La Tabla~\ref{tab:estado-arte} sintetiza los trabajos de referencia indicando su tarea, conjunto de datos, unidad de partición, validación externa, métricas y condiciones de comparabilidad.

\begin{sidewaystable}[p]
  \centering
  \caption{Comparación de trabajos relacionados en clasificación dermatológica.}
  \label{tab:estado-arte}
  \footnotesize
  \begin{tabularx}{\linewidth}{>{\hsize=0.85\hsize\linewidth=\hsize}Y >{\hsize=0.75\hsize\linewidth=\hsize}Y >{\hsize=0.95\hsize\linewidth=\hsize}Y >{\hsize=0.85\hsize\linewidth=\hsize}Y >{\hsize=0.85\hsize\linewidth=\hsize}Y >{\hsize=1.15\hsize\linewidth=\hsize}Y >{\hsize=1.60\hsize\linewidth=\hsize}Y}
    \toprule
    Trabajo & Tarea (n.º de clases) & Datos & Unidad de partición & Validación externa & Métrica reportada & Comparabilidad \\
    \midrule
    Hu et al.~\cite{hu2024multiscale} & Multiclase (7 clases) & HAM10000 sin vello e ISIC 2019 & No verificada & No verificada (evaluado en dos conjuntos) & Accuracy 94,0\,\% y AUC 99,3\,\% (HAM10000); accuracy 89,8\,\% (ISIC 2019) & Accuracy no equivale a Macro-F1; partición por lesión no verificada. \\
    Sarıateş y Özbay~\cite{sariates2025} & Clasificación dermatoscópica & Conjuntos públicos dermatoscópicos & No verificada & No verificada & Métricas de clasificación de ensamble & Protocolo de partición por lesión no verificado. \\
    Cherukuri y Vara Prasad~\cite{cherukuri2025} & Clasificación de cáncer de piel & Conjuntos dermatoscópicos públicos & No verificada & No verificada & Métricas de clasificación con aprendizaje por transferencia & Unidad de partición y desbalance no equiparables directamente. \\
    Barata et al.~\cite{barata2023reinforcement} & Multiclase (7 clases) & Entrenamiento en HAM10000; prueba en ISIC 2018 & No verificada en entrenamiento & Sí (ISIC 2018, 1.511 imágenes) & Sensibilidad MEL de 61,4\,\% a 79,5\,\% & 1.511 imágenes de prueba frente a 1.512; otro método de decisión (aprendizaje por refuerzo). \\
    Marchetti et al.~\cite{marchetti2023proveai} & Binaria (2 clases) & 603 lesiones de 435 participantes (95 melanomas) & Por paciente / prospectiva & Sí (estudio prospectivo PROVE-AI) & Sensibilidad 96,8\,\% y especificidad 37,4\,\% & Lesiones seleccionadas para biopsia; tarea binaria no comparable con 7 clases. \\
    Heinlein et al.~\cite{heinlein2024multicenter} & Binaria (2 clases) & 1.910 lesiones de 1.716 pacientes & Prospectiva multicéntrica & Sí (8 hospitales, 4 configuraciones de cámara) & IA: sensibilidad 92,1\,\% y especificidad 67,3\,\% & Lesiones sospechosas de melanoma en atención clínica; tarea binaria. \\
    Yan et al.~\cite{yan2025panderm} & Múltiple (28 tareas) & $>2$ millones de imágenes (11 instituciones, 4 modalidades) & Por tarea y cohorte & Sí (evaluación multicéntrica) & Desempeño en 28 tareas clínicas & Escala de datos y capacidad del modelo fundacional no comparables con MobileNetV3. \\
    Farías Rivera et al.~\cite{farias2026} & Clasificación en aplicativo web & Imágenes de lesiones cutáneas & No verificada & No verificada & Desempeño funcional y de clasificación & Antecedente de integración web; distinto protocolo experimental. \\
    \bottomrule
  \end{tabularx}
\end{sidewaystable}
\clearpage

Frente a los antecedentes resumidos en la Tabla~\ref{tab:estado-arte}, este trabajo integra cinco elementos metodológicos y de ingeniería en un mismo flujo reproducible. En primer lugar, aplica una partición estricta agrupada por lesión sobre las 10.015 imágenes de HAM10000~\cite{tschandl2018ham10000} y cuantifica experimentalmente la diferencia frente a una partición por imagen~\cite{kapoor2023leakage,abhishek2025quality}. En segundo lugar, compara ocho configuraciones de redes neuronales en resoluciones de 224 y 384\,px con tres semillas aleatorias (42, 43 y 44) por configuración. En tercer lugar, aplica una regla de selección fijada sobre el conjunto de validación antes de abrir el conjunto de prueba, la cual incorpora el desempate por AUC de melanoma~\cite{maierhein2024metrics}. Finalmente, evalúa el modelo calibrado ($T = 3{,}2933$) con sus umbrales fijos en las 1.512 imágenes de ISIC 2018~\cite{codella2019isic2018} y lo despliega en un aplicativo web que ejecuta el preprocesamiento y la inferencia localmente en el navegador~\cite{ning2024onnxweb}.

\chapter{Datos y preprocesamiento}

Este capítulo describe el conjunto de datos HAM10000~\cite{tschandl2018ham10000}, la estructura de múltiples fotografías por lesión y el procedimiento de partición agrupada en entrenamiento, validación y prueba. También detalla las cuatro etapas del preprocesamiento de imágenes~\cite{lee1997dullrazor,telea2004inpainting,zuiderveld1994clahe}, las transformaciones de aumento en línea y la ponderación de clases aplicada durante el entrenamiento. Su finalidad es documentar la preparación de los datos correspondiente al primer objetivo específico del trabajo.

\section{Conjunto de datos HAM10000}

El desarrollo experimental utilizó el conjunto público \textit{Human Against Machine with 10000 training images} (HAM10000), publicado por Tschandl et al.~\cite{tschandl2018ham10000} bajo licencia Creative Commons Attribution-NonCommercial (CC BY-NC). Este banco reúne 10.015 imágenes dermatoscópicas adquiridas durante 20 años en la Universidad Médica de Viena (Austria) y en la clínica de Cliff Rosendahl en Queensland (Australia)~\cite{tschandl2018ham10000}. Las 10.015 fotografías corresponden a 7.470 lesiones pigmentadas distribuidas en siete categorías diagnósticas~\cite{tschandl2018ham10000}. Más del 50\,\% de las lesiones cuenta con confirmación mediante estudio histopatológico, mientras que el resto se validó por seguimiento digital, microscopía confocal o consenso de expertos~\cite{tschandl2018ham10000}.

La Tabla~\ref{tab:distribucion} presenta el conteo y el porcentaje de imágenes por cada una de las siete clases del conjunto HAM10000~\cite{tschandl2018ham10000}. En ella se observa que la clase \texttt{nv} concentra 6.705 de las 10.015 fotografías, lo que equivale al 67,0\,\% del total del banco. En el extremo opuesto, las categorías con menor representación son \texttt{df}, con 115 imágenes (1,1\,\%), y \texttt{vasc}, con 142 imágenes (1,4\,\%)~\cite{tschandl2018ham10000}. La Figura~\ref{fig:clases} reserva el espacio para ilustrar un ejemplo dermatoscópico por cada una de las siete clases, atribuido a la colección original de HAM10000~\cite{tschandl2018ham10000}.

\begin{table}[H]
  \centering
  \caption{Distribución de imágenes por clase en HAM10000~\cite{tschandl2018ham10000}.}
  \label{tab:distribucion}
  \begin{tabularx}{\textwidth}{Yrr}
    \toprule
    Clase & Imágenes & Proporción (\%) \\
    \midrule
    \texttt{akiec} & 327 & 3,3 \\
    \texttt{bcc} & 514 & 5,1 \\
    \texttt{bkl} & 1.099 & 11,0 \\
    \texttt{df} & 115 & 1,1 \\
    \texttt{mel} & 1.113 & 11,1 \\
    \texttt{nv} & 6.705 & 67,0 \\
    \texttt{vasc} & 142 & 1,4 \\
    Total & 10.015 & 100,0 \\
    \bottomrule
  \end{tabularx}
\end{table}

\begin{figure}[H]
  \centering
  \fbox{\parbox[c][3cm][c]{0.85\linewidth}{\centering \pendiente{figura}}}
  \caption{Ejemplos dermatoscópicos de las siete clases de HAM10000~\cite{tschandl2018ham10000}.}
  \label{fig:clases}
\end{figure}

\section{Imágenes por lesión}

En los metadatos de HAM10000, cada imagen cuenta con un identificador de lesión (\texttt{lesion\_id}) que agrupa fotografías tomadas de la misma lesión en distintos aumentos, ángulos o momentos de seguimiento~\cite{tschandl2018ham10000}. La Tabla~\ref{tab:fotos-lesion} desglosa las 7.470 lesiones del conjunto según la cantidad de fotografías registradas por cada identificador. De las 7.470 lesiones, 5.514 cuentan con una única fotografía y 1.956 lesiones poseen entre dos y seis imágenes. Estas 1.956 lesiones con fotografías múltiples agrupan un total de 4.501 imágenes, lo que representa el 44,9\,\% de las 10.015 imágenes del conjunto.

\begin{table}[H]
  \centering
  \caption{Número de lesiones según la cantidad de fotos por lesión en HAM10000.}
  \label{tab:fotos-lesion}
  \begin{tabularx}{\textwidth}{Yr}
    \toprule
    Fotos por lesión & Número de lesiones \\
    \midrule
    1 & 5.514 \\
    2 & 1.423 \\
    3 & 490 \\
    4 & 34 \\
    5 & 5 \\
    6 & 4 \\
    Total & 7.470 \\
    \bottomrule
  \end{tabularx}
\end{table}

La repetición de fotografías por lesión no se distribuye de manera uniforme entre las categorías diagnósticas. En particular, la clase \texttt{mel} presenta un promedio de 1,81 fotografías por lesión, mientras que la clase mayoritaria \texttt{nv} registra un promedio de 1,24 fotografías por lesión. Esta diferencia implica que una división aleatoria por imagen tiene mayor probabilidad de asignar imágenes correlacionadas de un mismo melanoma tanto al entrenamiento como a la prueba~\cite{abhishek2025quality}. Por este motivo, el identificador \texttt{lesion\_id} se adoptó como unidad indivisible de agrupación al construir las particiones del estudio~\cite{kapoor2023leakage,abhishek2025quality}.

\section{Partición de los datos}

Para evitar que fotografías de una misma lesión quedaran repartidas entre distintos subconjuntos, la partición se realizó mediante el esquema \texttt{StratifiedGroupKFold} de la biblioteca scikit-learn~\cite{pedregosa2011sklearn} con semilla fija 42. El procedimiento se ejecutó en dos pasos sucesivos agrupando por \texttt{lesion\_id} y estratificando por la etiqueta de clase diagnóstica. En el primer paso, las 10.015 imágenes se dividieron en siete pliegues agrupados y se reservó un pliegue ($1/7$ del total) como conjunto de prueba. En el segundo paso, los seis pliegues restantes se subdividieron en seis grupos bajo el mismo criterio: uno ($1/7$ del total) pasó a validación y cinco ($5/7$) a entrenamiento.

Con esta división en séptimos, quedaron 7.153 imágenes (5.336 lesiones) en entrenamiento, 1.431 imágenes (1.067 lesiones) en validación y 1.431 imágenes (1.067 lesiones) en prueba. La Tabla~\ref{tab:particion} muestra la distribución exacta de imágenes por clase y el total de lesiones en cada uno de los tres subconjuntos. En los tres subconjuntos, cada una de las siete clases conserva su proporción original con diferencias menores o iguales a 0,1 puntos porcentuales. No se aplicó sobremuestreo ni submuestreo al construir las particiones para que los conjuntos de validación y prueba reflejaran la distribución original de HAM10000 sin alterar las prevalencias de evaluación~\cite{maierhein2024metrics,kapoor2023leakage}.

\begin{table}[H]
  \centering
  \caption{Partición agrupada por lesión según clase diagnóstica y número de lesiones.}
  \label{tab:particion}
  \begin{tabularx}{\textwidth}{Yrrr}
    \toprule
    Clase / unidad & Entrenamiento & Validación & Prueba \\
    \midrule
    \texttt{akiec} & 234 & 47 & 46 \\
    \texttt{bcc} & 366 & 74 & 74 \\
    \texttt{bkl} & 785 & 157 & 157 \\
    \texttt{df} & 82 & 16 & 17 \\
    \texttt{mel} & 795 & 159 & 159 \\
    \texttt{nv} & 4.789 & 958 & 958 \\
    \texttt{vasc} & 102 & 20 & 20 \\
    Total imágenes & 7.153 & 1.431 & 1.431 \\
    Total lesiones & 5.336 & 1.067 & 1.067 \\
    \bottomrule
  \end{tabularx}
\end{table}

Dado que Abhishek et al.~\cite{abhishek2025quality} identificaron 18 pares de imágenes duplicadas en HAM10000 que poseen identificadores \texttt{lesion\_id} distintos, se verificó la ubicación de dichos pares dentro de la partición generada. De los 18 pares reportados~\cite{abhishek2025quality}, 14 pares quedaron asignados dentro del mismo subconjunto y cuatro pares quedaron separados entre particiones. En concreto, dos pares quedaron entre prueba y entrenamiento (ambos \texttt{bkl}) y dos entre validación y entrenamiento (ambos \texttt{nv}), afectando a 2 de las 1.431 imágenes de prueba (0,14\,\%). Se deja constancia de que agrupar por \texttt{lesion\_id} no garantiza separación por paciente, puesto que HAM10000 no publica identificadores de paciente en sus metadatos~\cite{tschandl2018ham10000,abhishek2025quality}.

\section{Preprocesamiento}

Todas las imágenes de entrenamiento, validación y prueba se sometieron a una cadena determinista de cuatro etapas antes de ingresar a las redes neuronales. En la primera etapa, cada fotografía se redimensionó a 224\,px o 384\,px conservando su proporción original mediante relleno simétrico de píxeles negros en los bordes. Este relleno evita la distorsión geométrica de la asimetría de la lesión que produciría un escalado anisotrópico directo. A continuación, en la segunda etapa se aplicó un filtro de mediana con ventana de $3 \times 3$ píxeles para atenuar el ruido impulsivo sin difuminar los bordes de la lesión.

En la tercera etapa se implementó un módulo de eliminación de vello basado en operaciones morfológicas e interpolación~\cite{lee1997dullrazor,telea2004inpainting}. Siguiendo el método morfológico de DullRazor~\cite{lee1997dullrazor}, se aplicó una transformada \textit{black-hat} en escala de grises y un umbral binario para aislar las hebras oscuras de vello. Los píxeles de la máscara se reconstruyeron mediante el \textit{inpainting} de marcha rápida de Telea~\cite{telea2004inpainting}, que propaga el color desde el contorno hacia el interior ocluido. En HAM10000, la oclusión abundante de vello reduce el rendimiento (afectando más a \texttt{vasc} y menos a \texttt{akiec}), aunque el vello escaso a veces no perjudica la predicción~\cite{wang2024hair}.

En la cuarta etapa se aplicó la ecualización adaptativa de histograma limitada por contraste (CLAHE)~\cite{zuiderveld1994clahe} con un límite de recorte (\textit{clip limit}) de 2,0. Esta operación realza el contraste local de las redes de pigmento y las estructuras vasculares a la vez que acota la amplificación del ruido en regiones homogéneas de piel sana~\cite{zuiderveld1994clahe}. Se aclara que los parámetros de estas cuatro etapas se fijaron en la cadena base y no se realizó un estudio de ablación por etapa. La Figura~\ref{fig:preprocesamiento} reserva el espacio para mostrar el resultado visual antes y después de aplicar la secuencia de preprocesamiento sobre una imagen dermatoscópica.

\begin{figure}[H]
  \centering
  \fbox{\parbox[c][3cm][c]{0.85\linewidth}{\centering \pendiente{figura}}}
  \caption{Ilustración del antes y después de las etapas de preprocesamiento dermatoscópico.}
  \label{fig:preprocesamiento}
\end{figure}

\section{Aumento de datos}

Para incrementar la variabilidad geométrica y fotométrica sin alterar los conjuntos de evaluación, se aplicó aumento de datos en línea exclusivamente sobre las 7.153 imágenes de entrenamiento~\cite{zhang2024dive}. En cada época se aplicaron volteos horizontal y vertical, rotación de $\pm 20^\circ$, zoom de $\pm 10$\,\%, brillo y contraste de $\pm 0{,}15$ y relleno exterior por reflexión. No se aplicaron transformaciones de matiz ni saturación de tono (\textit{hue}) porque la coloración rojiza, azulada o marrón oscura constituye un rasgo diagnóstico entre lesiones vasculares, melanomas y nevos~\cite{tschandl2018ham10000}.

Antes de adoptar el esquema en línea, se evaluó y descartó una versión previa de aumento estático en disco que generaba 98.239 imágenes almacenadas antes del entrenamiento. En las pruebas preliminares, ese volumen de archivos elevaba el tiempo de cómputo a cerca de 8 minutos por época debido al cuello de botella de lectura en disco. Además, el sobremuestreo físico en esas 98.239 imágenes alteraba las frecuencias de clase y anulaba el efecto de los pesos de clase en la función de pérdida.

\section{Tratamiento del desbalance}

Para compensar la desproporción entre las 4.789 imágenes de \texttt{nv} y las 82 de \texttt{df} en entrenamiento (Tabla~\ref{tab:particion}), se incorporaron pesos de clase en la entropía cruzada~\cite{hu2024multiscale}. En lugar de emplear el inverso directo de la frecuencia, se definió una ponderación suavizada mediante la raíz cuadrada del inverso de la proporción de cada clase $c$. Con $N = 7.153$ imágenes totales de entrenamiento y $n_c$ imágenes en la clase $c$, el peso $w_c$ se calculó mediante la Ecuación~\ref{eq:pesos-clase}:

\begin{equation}
  w_c = \sqrt{\frac{N}{7 \cdot n_c}}
  \label{eq:pesos-clase}
\end{equation}

La Tabla~\ref{tab:pesos} presenta los pesos resultantes para cada una de las siete clases diagnósticas en el conjunto de entrenamiento. Si se hubiera utilizado la ponderación inversa directa $N / (7 \cdot n_c)$, la clase \texttt{df} ($n_{\texttt{df}} = 82$) habría recibido un peso 58 veces mayor que el de la clase \texttt{nv} ($n_{\texttt{nv}} = 4.789$). Con la formulación de raíz cuadrada de la Ecuación~\ref{eq:pesos-clase}, la relación entre el peso máximo (\texttt{df}, con 3,53) y el mínimo (\texttt{nv}, con 0,46) se acota a unas 7,7 veces. Mientras otros trabajos combinan ponderación con remuestreo o \textit{label smoothing}~\cite{hu2024multiscale}, esta raíz cuadrada evitó la inestabilidad de gradiente observada en las corridas de diagnóstico.

\begin{table}[H]
  \centering
  \caption{Pesos de clase empleados en la función de pérdida durante el entrenamiento ($N = 7.153$).}
  \label{tab:pesos}
  \begin{tabularx}{\textwidth}{Yrr}
    \toprule
    Clase & Imágenes en entrenamiento ($n_c$) & Peso ($w_c$) \\
    \midrule
    \texttt{nv} & 4.789 & 0,46 \\
    \texttt{mel} & 795 & 1,13 \\
    \texttt{bkl} & 785 & 1,14 \\
    \texttt{bcc} & 366 & 1,67 \\
    \texttt{akiec} & 234 & 2,09 \\
    \texttt{vasc} & 102 & 3,17 \\
    \texttt{df} & 82 & 3,53 \\
    \bottomrule
  \end{tabularx}
\end{table}

\chapter{Modelos y entrenamiento}

Este capítulo documenta las siete arquitecturas evaluadas, la receta de entrenamiento en dos fases y los ajustes derivados del diagnóstico inicial en validación~\cite{zhang2024dive}. Además, presenta la comparación experimental entre la partición por imagen y por lesión~\cite{kapoor2023leakage}, el escalado a 384\,px y la regla preregistrada que seleccionó MobileNetV3Large~\cite{howard2019mobilenetv3}. Por último, describe el ajuste de calibración por temperatura ($T = 3{,}2933$)~\cite{guo2017calibration} y la definición de los umbrales de decisión para melanoma~\cite{barata2023reinforcement}.

\section{Arquitecturas evaluadas}

Para cubrir el segundo objetivo específico, se implementaron siete arquitecturas de aprendizaje profundo: cinco bases preentrenadas en ImageNet~\cite{deng2009imagenet}, una variante híbrida preentrenada y un modelo híbrido entrenado desde cero. Las cinco redes convolucionales puras correspondieron a DenseNet121~\cite{huang2017densenet}, ResNet50~\cite{he2016resnet}, MobileNetV3Large~\cite{howard2019mobilenetv3}, EfficientNetV2-S~\cite{tan2021efficientnetv2} y ConvNeXt-Tiny~\cite{liu2022convnext}. Las dos arquitecturas con capas recurrentes fueron DenseNet121-BiLSTM (con base preentrenada en ImageNet) y una red CNN-BiLSTM entrenada íntegramente desde cero sobre HAM10000~\cite{zhang2024dive,tschandl2018ham10000}.

La Tabla~\ref{tab:arquitecturas} resume cada arquitectura junto con su número de parámetros, el origen de sus pesos iniciales y las resoluciones espaciales en las que fue evaluada. En ella se observa que MobileNetV3Large presenta el menor tamaño entre las redes caracterizadas a 384\,px, con 3,24 millones de parámetros, mientras que ConvNeXt-Tiny alcanza 28,02 millones de parámetros~\cite{howard2019mobilenetv3,liu2022convnext}. En MobileNetV3Large, la cabeza clasificadora incluyó agrupamiento promedio global, \textit{Dropout} de 0,3, capa densa de 256 unidades con ReLU, \textit{Dropout} de 0,2 y salida de siete unidades con \textit{softmax}~\cite{howard2019mobilenetv3,zhang2024dive}. La Figura~\ref{fig:arquitectura-modelo} reserva el espacio para el diagrama de bloques de esta arquitectura seleccionada.

\begin{table}[H]
  \centering
  \caption{Arquitecturas evaluadas, número de parámetros, preentrenamiento y resoluciones de entrada.}
  \label{tab:arquitecturas}
  \begin{tabularx}{\textwidth}{>{\hsize=1.3\hsize\linewidth=\hsize}Y >{\hsize=0.8\hsize\linewidth=\hsize}Y >{\hsize=0.95\hsize\linewidth=\hsize}Y >{\hsize=0.95\hsize\linewidth=\hsize}Y}
    \toprule
    Arquitectura & Parámetros (M) & Preentrenamiento & Resoluciones evaluadas \\
    \midrule
    MobileNetV3Large~\cite{howard2019mobilenetv3} & 3,24 & ImageNet~\cite{deng2009imagenet} & 224\,px y 384\,px \\
    DenseNet121~\cite{huang2017densenet} & 7,30 & ImageNet~\cite{deng2009imagenet} & 224\,px y 384\,px \\
    EfficientNetV2-S~\cite{tan2021efficientnetv2} & 20,66 & ImageNet~\cite{deng2009imagenet} & 224\,px y 384\,px \\
    ConvNeXt-Tiny~\cite{liu2022convnext} & 28,02 & ImageNet~\cite{deng2009imagenet} & 384\,px (cribado) \\
    ResNet50~\cite{he2016resnet} & \pendiente{parámetros ResNet50} & ImageNet~\cite{deng2009imagenet} & 224\,px \\
    DenseNet121-BiLSTM~\cite{huang2017densenet,zhang2024dive} & \pendiente{parámetros DenseNet121-BiLSTM} & ImageNet (base)~\cite{deng2009imagenet} & 224\,px \\
    CNN-BiLSTM (desde cero)~\cite{zhang2024dive} & \pendiente{parámetros CNN-BiLSTM} & Ninguno (desde cero) & 224\,px \\
    \bottomrule
  \end{tabularx}
\end{table}

\begin{figure}[H]
  \centering
  \fbox{\parbox[c][3cm][c]{0.85\linewidth}{\centering \pendiente{figura}}}
  \caption{Diagrama de la arquitectura MobileNetV3Large y su cabeza de clasificación.}
  \label{fig:arquitectura-modelo}
\end{figure}

\section{Estrategia de entrenamiento}

Todas las redes preentrenadas se ajustaron en dos fases con el optimizador Adam, entropía cruzada con pesos de clase (Tabla~\ref{tab:pesos}) y precisión mixta~\cite{zhang2024dive}. En la Fase~1 se congeló la base convolucional y se entrenó solo la cabeza clasificadora con tasa de $3 \cdot 10^{-4}$ hasta por 15 épocas. En la Fase~2 se descongeló toda la red excepto las capas de normalización por lotes (BatchNorm) durante un máximo de 40 épocas. En esta segunda etapa se aplicó un calentamiento lineal (\textit{warmup}) de 2 épocas desde $1 \cdot 10^{-5}$ hasta $5 \cdot 10^{-5}$, seguido de un decaimiento coseno hasta $1 \cdot 10^{-6}$.

Las capas BatchNorm permanecieron congeladas en la Fase~2 para conservar las estadísticas móviles de ImageNet~\cite{deng2009imagenet} y evitar oscilaciones con lotes de 32 o 64 imágenes~\cite{zhang2024dive}. La parada temprana y la reducción de tasa se guiaron por el Macro-F1 de validación, con paciencia de 15 épocas y un mínimo de 20 épocas. Por su parte, la red CNN-BiLSTM se entrenó desde cero en una fase de hasta 100 épocas con lote de 64 y decaimiento coseno ($1 \cdot 10^{-5} \to 1 \cdot 10^{-3} \to 1 \cdot 10^{-5}$).

El lote fue de 64 imágenes a 224\,px y de 32 a 384\,px, ejecutándose en una estación con GPU NVIDIA RTX A4000 del laboratorio universitario. Cada configuración se entrenó con tres semillas aleatorias independientes (42, 43 y 44) y las evaluaciones utilizaron los pesos de la última época sin aumento en tiempo de prueba (TTA). El conjunto de prueba se usó dos veces (fase de 224\,px y fase de 384\,px), siempre con la selección fijada antes en validación~\cite{kapoor2023leakage}.

\section{Diagnóstico inicial y ajustes}

El punto de partida experimental con DenseNet121 a 224\,px produjo en su primera corrida un Macro-F1 de validación de 0,618. Para identificar los factores que limitaban el aprendizaje de las clases minoritarias, se ejecutaron seis corridas sucesivas de diagnóstico sobre el conjunto de validación. El diagnóstico reveló tres causas: caída prematura de la tasa de aprendizaje, pesos inversos directos excesivos (\texttt{df} pesaba 58 veces más que \texttt{nv}) y ajuste fino insuficiente.

El calentamiento con decaimiento coseno, los pesos por raíz cuadrada (Ecuación~\ref{eq:pesos-clase}) y la ampliación de la Fase~2 a 40 épocas corrigieron estas limitaciones. Tras aplicar estos tres ajustes en la receta base, DenseNet121 a 224\,px elevó su Macro-F1 de validación de 0,618 a 0,685 y alcanzó una accuracy de validación de 0,806. Esta receta estabilizada se fijó como protocolo estándar para comparar todas las arquitecturas siguientes.

\section{Comparación a 224 px y ronda de regularización}

En la evaluación comparativa a 224\,px sobre el conjunto de validación, todas las arquitecturas con preentrenamiento en ImageNet~\cite{deng2009imagenet} se ubicaron en el rango de 0,66 a 0,70 de Macro-F1 medio. En contraste, la red CNN-BiLSTM entrenada desde cero alcanzó un Macro-F1 de validación de 0,56. Como hipótesis, una explicación posible es que 7.153 imágenes no bastan para aprender desde cero los filtros convolucionales y las capas recurrentes sin preentrenamiento~\cite{zhang2024dive}.

En MobileNetV3Large a 224\,px se obtuvo un Macro-F1 medio de validación de 0,697 y una brecha de $+0{,}24$ puntos entre entrenamiento y validación. Con el objetivo de reducir esa brecha de generalización, se ejecutó una ronda específica de regularización evaluada en validación con tres semillas por variante. La primera variante incorporó \textit{label smoothing} de 0,1~\cite{hu2024multiscale} y el optimizador AdamW, con lo cual el Macro-F1 medio de validación descendió de 0,697 a 0,680. La segunda variante probó aumento fuerte (rotación $\pm 30^\circ$, zoom $\pm 20$\,\%, brillo/contraste $\pm 0{,}2$ y borrado del 5--20\,\% del área), bajando el Macro-F1 a 0,625; ninguna variante se adoptó.

\section{Partición por imagen frente a partición por lesión}

Para cuantificar el efecto de la unidad de partición sobre las métricas de evaluación~\cite{kapoor2023leakage,abhishek2025quality}, se construyó en paralelo una partición aleatoria por imagen a 224\,px. Esta división utilizó la función \texttt{train\_test\_split} de scikit-learn~\cite{pedregosa2011sklearn} en proporción estratificada 70/15/15 con semilla 42, lo que generó 7.010 imágenes de entrenamiento, 1.502 de validación y 1.503 de prueba. Al cruzar los \texttt{lesion\_id}, se comprobó que el 36,1\,\% de la prueba (543 de 1.503 imágenes) tenía otra fotografía de la misma lesión en entrenamiento. En el caso específico de la clase \texttt{mel}, el 63,5\,\% de los melanomas de prueba (106 de 167) tenía otra imagen de la misma lesión dentro del entrenamiento.

La Tabla~\ref{tab:comparacion-particiones} compara la media de tres semillas en prueba a 224\,px para DenseNet121 y CNN-BiLSTM bajo ambos protocolos de partición. En DenseNet121, partir por imagen subió la accuracy en $+3{,}3$ puntos porcentuales y el Macro-F1 en $+8{,}7$ puntos (de 62,9\,\% a 71,6\,\%). En la arquitectura CNN-BiLSTM, la partición por imagen aumentó la accuracy en $+4{,}2$ puntos y el Macro-F1 en $+9{,}2$ puntos porcentuales (de 49,2\,\% a 58,4\,\%).

\begin{table}[H]
  \centering
  \caption{Comparación entre partición por imagen y partición por lesión a 224\,px (prueba, media de tres semillas).}
  \label{tab:comparacion-particiones}
  \begin{tabularx}{\textwidth}{YYrr}
    \toprule
    Modelo & Unidad de partición & Accuracy & Macro-F1 \\
    \midrule
    DenseNet121 & Por imagen & 82,8\,\% & 71,6\,\% \\
    DenseNet121 & Por lesión & 79,5\,\% & 62,9\,\% \\
    CNN-BiLSTM & Por imagen & 76,8\,\% & 58,4\,\% \\
    CNN-BiLSTM & Por lesión & 72,6\,\% & 49,2\,\% \\
    \bottomrule
  \end{tabularx}
\end{table}

Estos incrementos de $+8{,}7$ a $+9{,}2$ puntos en Macro-F1 aportan evidencia empírica de inflación de métricas por fuga de información entre entrenamiento y prueba~\cite{kapoor2023leakage,abhishek2025quality}. Se aclara que ambos conjuntos de prueba no contienen las mismas imágenes (1.503 frente a 1.431), por lo que no toda la diferencia se atribuye a la fuga~\cite{kapoor2023leakage}. En virtud de la contaminación comprobada en el 36,1\,\% de la prueba, la partición por imagen se descartó definitivamente para el resto del estudio.

\section{Resolución de 384 px}

Para evaluar el efecto de una mayor resolución en la partición por lesión, se ejecutó un cribado a 384\,px en validación (semilla 42) con cuatro arquitecturas. La Tabla~\ref{tab:cribado-384} presenta el Macro-F1 de validación, el conteo de parámetros y el tiempo promedio por época de cada candidato. Las tres primeras arquitecturas superaron el umbral de 0,73 de Macro-F1 en validación, encabezadas por MobileNetV3Large con 0,757 y 28,3\,s por época.

\begin{table}[H]
  \centering
  \caption{Cribado de arquitecturas a 384\,px en el conjunto de validación (semilla 42).}
  \label{tab:cribado-384}
  \begin{tabularx}{\textwidth}{Yrrr}
    \toprule
    Arquitectura & Macro-F1 (validación) & Parámetros (M) & Tiempo por época (s) \\
    \midrule
    MobileNetV3Large~\cite{howard2019mobilenetv3} & 0,757 & 3,24 & 28,3 \\
    EfficientNetV2-S~\cite{tan2021efficientnetv2} & 0,749 & 20,66 & 34,7 \\
    DenseNet121~\cite{huang2017densenet} & 0,733 & 7,30 & 35,3 \\
    ConvNeXt-Tiny~\cite{liu2022convnext} & 0,729 & 28,02 & 37,0 \\
    \bottomrule
  \end{tabularx}
\end{table}

Según la Tabla~\ref{tab:cribado-384}, ConvNeXt-Tiny~\cite{liu2022convnext} se descartó tras el cribado y no se entrenó con las semillas 43 y 44 ni se evaluó en prueba. Esta exclusión se debió a que obtuvo el menor Macro-F1 en validación (0,729), el mayor número de parámetros (28,02\,M) y el mayor tiempo por época (37,0\,s). Como hipótesis, una explicación posible es que la resolución de 384\,px conserva detalles morfológicos finos de la lesión, como bordes irregulares y glóbulos de pigmento.

\section{Regla de selección del modelo}

Antes de abrir la prueba a 384\,px, se fijó una regla de selección basada únicamente en las métricas del conjunto de validación~\cite{kapoor2023leakage,maierhein2024metrics}. En el primer paso se calcula para cada arquitectura la media del Macro-F1 de validación entre las tres semillas (42, 43 y 44). En el segundo paso empatan las arquitecturas situadas a menos de $\delta = 0{,}0112$ de la mejor media, donde $\delta$ es la raíz de la varianza media entre semillas. En el tercer paso desempatan por AUC media de melanoma en validación; si difieren en menos de 0,005, gana la de menos parámetros y se despliega la semilla mediana.

Al finalizar las tres semillas a 384\,px en validación, EfficientNetV2-S obtuvo la mayor media de Macro-F1 con 0,748, seguida por MobileNetV3Large con 0,738. Dado que la diferencia entre ambas medias fue de 0,0100 (menor que el margen de empate $\delta = 0{,}0112$), las dos arquitecturas empataron en el primer criterio. En el desempate por AUC de melanoma en validación, MobileNetV3Large superó a EfficientNetV2-S con 0,901 frente a 0,886 (diferencia de 0,015, mayor que 0,005).

Se deja constancia de que, si la selección se hubiera basado únicamente en el promedio de Macro-F1 sin margen de empate, habría resultado elegida EfficientNetV2-S (0,748 frente a 0,738). El desempate por AUC de melanoma responde a la relevancia clínica de esta neoplasia~\cite{amaral2025esmo,maierhein2024metrics} y seleccionó un modelo con 6,4 veces menos parámetros (3,24\,M frente a 20,66\,M). De las tres corridas de MobileNetV3Large a 384\,px, la regla seleccionó para el despliegue y la validación externa la semilla 44, correspondiente al valor mediano de Macro-F1 en validación.

\section{Calibración y umbrales de decisión}

Una vez seleccionado el modelo MobileNetV3Large a 384\,px (semilla 44), se evaluó la calibración de sus probabilidades sobre las 1.431 imágenes del conjunto de validación~\cite{sambyal2023calibration,penso2024calibration}. En clasificación de imágenes médicas, las redes profundas entrenadas con entropía cruzada ponderada suelen producir distribuciones \textit{softmax} concentradas en valores extremos que no coinciden con la probabilidad empírica de acierto~\cite{guo2017calibration,sambyal2023calibration}. Para corregir esta sobreconfianza antes de fijar puntos de corte clínicos, se ajustó sobre el conjunto de validación el parámetro de \textit{temperature scaling} (Ecuación~\ref{eq:tempscaling})~\cite{guo2017calibration,penso2024calibration}, obteniéndose una temperatura de $T = 3{,}2933$.

Con los \textit{logits} escalados por $T = 3{,}2933$, se definieron modos que sustituyen la clase más probable (\textit{argmax}) por un umbral $t$ sobre $p_{\text{cal}}(\texttt{mel})$. Bajo esta regla jerárquica, el sistema predice la clase \texttt{mel} siempre que se cumpla la condición de la Ecuación~\ref{eq:regla-umbral}:

\begin{equation}
  \hat{y} =
  \begin{cases}
    \texttt{mel}, & \text{si } p_{\text{cal}}(\texttt{mel}) \ge t \\
    \argmax_{c \neq \texttt{mel}} p_{\text{cal}}(c), & \text{si } p_{\text{cal}}(\texttt{mel}) < t
  \end{cases}
  \label{eq:regla-umbral}
\end{equation}

Sobre la curva de validación del modelo calibrado se establecieron dos umbrales operativos para el aplicativo web: un modo equilibrado con $t = 0{,}20$ y un modo de alta sensibilidad con $t = 0{,}11$. El propósito de ambos modos es intercambiar falsos negativos de melanoma por falsos positivos, dado que omitir una lesión maligna conlleva un riesgo clínico mayor que derivar un nevo benigno a revisión dermatológica~\cite{barata2023reinforcement,amaral2025esmo}. Esta configuración recoge también la evidencia del estudio prospectivo PROVE-AI~\cite{marchetti2023proveai}, según la cual elevar la sensibilidad mediante umbrales bajos incrementa la tasa de falsos positivos y exige reportar conjuntamente sensibilidad y especificidad.

\chapter{Construcción de la aplicación}

Este capítulo documenta la arquitectura, los requisitos, el diseño de interfaz y las funcionalidades del aplicativo web desarrollado para el cuarto objetivo específico. Describe la exportación del modelo MobileNetV3Large a formato ONNX y su ejecución en el navegador mediante OpenCV.js y ONNX Runtime Web~\cite{ning2024onnxweb}. Por último, presenta el manual de usuario y los resultados de verificación técnica en paridad numérica, rendimiento, memoria~\cite{wang2025browser} y accesibilidad automática~\cite{w3c2023wcag22}.

\section{Arquitectura del sistema}

La arquitectura del aplicativo se diseñó bajo un esquema de ejecución en el cliente (\textit{client-side}) en el que todas las etapas de procesamiento ocurren dentro del navegador del usuario~\cite{ning2024onnxweb}. Esta decisión arquitectónica garantiza que las imágenes dermatoscópicas no salgan del equipo local, elimina la necesidad de servidores dedicados de inferencia y suprime los costos de cómputo en la nube~\cite{ning2024onnxweb}. En las siguientes subsecciones se detallan la visión general, los módulos del cliente, la exportación a ONNX, el flujo de datos y la configuración de despliegue estático.

\subsection{Visión general}

El aplicativo web se construyó con React, TypeScript y el empaquetador Vite como una aplicación de página única servida de manera estática. Para ejecutar el análisis sin un servidor de cómputo, el cliente integra dos motores en WebAssembly: OpenCV.js 5.0.0 y \texttt{onnxruntime-web} 1.31.1 con cuatro hilos~\cite{ning2024onnxweb}. Tras descargar los archivos estáticos y el modelo ONNX de 12,9\,MB (1,5\,s en Chromium 158), todo el procesamiento ocurre en la memoria local del navegador.

El diseño arquitectónico del sistema se representa mediante dos vistas complementarias. La Figura~\ref{fig:contexto} reserva el espacio para el diagrama de contexto del aplicativo, donde el usuario analiza imágenes sin enviarlas hacia servicios externos. Por su parte, la Figura~\ref{fig:c4-contenedores} reserva el espacio para el diagrama C4 de contenedores, que separa la interfaz en el hilo principal del cómputo en un \textit{Web Worker}~\cite{wang2025browser,ning2024onnxweb}.

\begin{figure}[H]
  \centering
  \fbox{\parbox[c][3cm][c]{0.85\linewidth}{\centering \pendiente{figura}}}
  \caption{Diagrama de contexto del aplicativo web con inferencia local.}
  \label{fig:contexto}
\end{figure}

\begin{figure}[H]
  \centering
  \fbox{\parbox[c][3cm][c]{0.85\linewidth}{\centering \pendiente{figura}}}
  \caption{Diagrama C4 de contenedores de la arquitectura en el navegador.}
  \label{fig:c4-contenedores}
\end{figure}

\subsection{Módulos del cliente}

El código del cliente se divide en dos contextos de ejecución que se comunican mediante paso de mensajes asíncronos. El hilo principal en React y TypeScript gestiona la carga de archivos (individuales, múltiples, carpetas o ZIP) y el selector de modo ($t = 0{,}20$ o $t = 0{,}11$). También presenta la tabla de resultados, el detalle por imagen y la exportación a valores separados por comas (CSV), notación de objetos de JavaScript (JSON) e impresión. Este hilo controla el estado visual de la aplicación y valida el formato de los archivos antes de encolarlos para su análisis.

El segundo contexto corresponde a un \textit{Web Worker} dedicado que encapsula la instancia de OpenCV.js 5.0.0 y la sesión de inferencia de \texttt{onnxruntime-web} 1.31.1~\cite{ning2024onnxweb}. Esta separación responde a que ejecutar la inferencia en el hilo principal del navegador aumenta en promedio un 67,2\,\% el tiempo de renderizado de la interfaz~\cite{wang2025browser}. Al aislar el preprocesamiento y la evaluación del modelo en el \textit{Web Worker}, la interfaz mantiene su capacidad de respuesta mientras se procesan lotes de hasta 300 imágenes consecutivas.

\subsection{Exportación del modelo a ONNX}

Para ejecutar el clasificador seleccionado (MobileNetV3Large a 384\,px, semilla 44)~\cite{howard2019mobilenetv3} en el navegador, el modelo de Keras se exportó a ONNX con \textit{opset} 17~\cite{ning2024onnxweb}. Como el entrenamiento usó precisión mixta, los pesos y tensores se convirtieron a punto flotante de 32 bits (\texttt{float32}) para generar el archivo de 12,9\,MB. El grafo exportado recibe un tensor en formato de lote, altura, anchura y canales de $1 \times 384 \times 384 \times 3$ y entrega los \textit{logits} de las siete clases diagnósticas.

Debido a la conversión a \texttt{float32}, se midió la diferencia entre el modelo original en Keras y el archivo ONNX ejecutado en el mismo equipo. En esta comprobación, la probabilidad de melanoma mostró una diferencia en el percentil 99 (p99) de 0,0236 y una diferencia absoluta máxima de 0,0587 (\pendiente{conjunto de imágenes sobre el que se midió y archivo del informe}). Esta cuantificación separa el efecto de la conversión a \texttt{float32} de la verificación posterior entre los motores ONNX de Python y del navegador.

\subsection{Flujo de datos}

El aplicativo procesa cada lote de imágenes mediante una secuencia determinista entre el hilo principal y el \textit{Web Worker}. La Figura~\ref{fig:casos-uso} reserva el espacio para el diagrama de casos de uso que resume las operaciones disponibles para el usuario. De forma complementaria, la Figura~\ref{fig:secuencia} reserva el espacio para el diagrama de secuencia de los siete pasos numerados que se describen a continuación:

\begin{enumerate}
  \item \textbf{Carga y validación de entrada:} el usuario selecciona una imagen, múltiples archivos, una carpeta o un contenedor ZIP; el hilo principal extrae los archivos válidos y descarta aquellos con formato no soportado.
  \item \textbf{Decodificación de imagen:} cada archivo se decodifica a una matriz de píxeles en color (canales rojo, verde y azul) y se transfiere al \textit{Web Worker} para iniciar la tubería de análisis.
  \item \textbf{Preprocesamiento en OpenCV.js 5.0.0:} el \textit{Web Worker} aplica en orden el redimensionado a $384 \times 384$\,px con relleno negro que conserva la proporción, el filtro de mediana de $3 \times 3$, la eliminación de vello por transformada \textit{black-hat}, umbral e \textit{inpainting} de Telea~\cite{lee1997dullrazor,telea2004inpainting} y la ecualización CLAHE con límite de recorte de 2,0~\cite{zuiderveld1994clahe}.
  \item \textbf{Inferencia en \texttt{onnxruntime-web} 1.31.1:} la matriz preprocesada se empaqueta como tensor \texttt{float32} y se evalúa sobre el modelo ONNX (opset 17, 12,9\,MB) mediante el motor WebAssembly configurado con cuatro hilos~\cite{ning2024onnxweb}.
  \item \textbf{Calibración de probabilidades:} los siete \textit{logits} devueltos por la red se dividen por la temperatura $T = 3{,}2933$ antes de calcular la función \textit{softmax} (Ecuación~\ref{eq:tempscaling})~\cite{guo2017calibration,penso2024calibration}.
  \item \textbf{Regla de decisión por modo:} el sistema compara la probabilidad calibrada $p_{\text{cal}}(\texttt{mel})$ con el umbral del modo activo ($t = 0{,}20$ en modo equilibrado o $t = 0{,}11$ en modo de alta sensibilidad); si lo iguala o supera, asigna \texttt{mel}, y en caso contrario selecciona la clase más probable entre las otras seis categorías (Ecuación~\ref{eq:regla-umbral}).
  \item \textbf{Presentación y exportación:} el \textit{Web Worker} libera los tensores temporales de memoria y devuelve el resultado al hilo principal, que actualiza la tabla de resumen, habilita la vista de detalle por imagen y ofrece la descarga en CSV, JSON o vista de impresión.
\end{enumerate}

\begin{figure}[H]
  \centering
  \fbox{\parbox[c][3cm][c]{0.85\linewidth}{\centering \pendiente{figura}}}
  \caption{Diagrama de casos de uso del aplicativo web.}
  \label{fig:casos-uso}
\end{figure}

\begin{figure}[H]
  \centering
  \fbox{\parbox[c][3cm][c]{0.85\linewidth}{\centering \pendiente{figura}}}
  \caption{Diagrama de secuencia del flujo de preprocesamiento e inferencia en el navegador.}
  \label{fig:secuencia}
\end{figure}

\subsection{Despliegue en Vercel}

El aplicativo compilado mediante Vite se desplegó como sitio web estático en la plataforma Vercel en la dirección \pendiente{URL}. En este esquema, Vercel sirve únicamente archivos estáticos (HTML, hojas de estilo, JavaScript, binarios WebAssembly y el modelo ONNX de 12,9\,MB), sin ejecutar código en servidor~\cite{ning2024onnxweb}. Para habilitar cuatro hilos de WebAssembly con memoria compartida (\texttt{SharedArrayBuffer})~\cite{ning2024onnxweb}, se configuraron las cabeceras HTTP de aislamiento de origen (\texttt{Cross-Origin-Opener-Policy: same-origin} y \texttt{Cross-Origin-Embedder-Policy: require-corp}). Esta configuración habilita el paralelismo en CPU con los tiempos de respuesta medidos en la Sección~6.6 (Tabla~\ref{tab:requisitos-no-funcionales}).

\section{Especificación de requisitos}

La ingeniería del aplicativo se formalizó mediante seis requisitos funcionales (RF-01 a RF-06) y cinco requisitos no funcionales (RNF-01 a RNF-05) estructurados según el estándar IEEE. La Tabla~\ref{tab:requisitos-funcionales} detalla los seis requisitos funcionales con su identificador, descripción, prioridad, criterio de aceptación y estado de verificación documental. Como los informes del repositorio cubrieron paridad, rendimiento, memoria y accesibilidad, la columna de verificación de los RF deja el marcador donde no hay un acta separada.

\begin{table}[H]
  \centering
  \caption{Requisitos funcionales del aplicativo web.}
  \label{tab:requisitos-funcionales}
  \footnotesize
  \begin{tabularx}{\textwidth}{l >{\hsize=0.95\hsize\linewidth=\hsize}Y >{\hsize=0.55\hsize\linewidth=\hsize}Y >{\hsize=1.55\hsize\linewidth=\hsize}Y >{\hsize=0.95\hsize\linewidth=\hsize}Y}
    \toprule
    ID & Descripción & Prioridad & Criterio de aceptación & Verificación \\
    \midrule
    RF-01 & Cargar una imagen individual & \pendiente{prioridad} & El sistema acepta un archivo de imagen válido, ejecuta la cadena de 7 pasos y muestra su predicción. & \pendiente{} \\
    RF-02 & Cargar varias imágenes simultáneamente & \pendiente{prioridad} & El sistema procesa en lote una selección múltiple de imágenes e informa los archivos rechazados. & \pendiente{} \\
    RF-03 & Cargar una carpeta de imágenes & \pendiente{prioridad} & El sistema recorre los archivos de imagen de un directorio local y genera la tabla consolidada de resultados. & \pendiente{} \\
    RF-04 & Cargar un archivo comprimido ZIP & \pendiente{prioridad} & El sistema descomprime en memoria el archivo ZIP, filtra las imágenes válidas y ejecuta la inferencia por lote. & \pendiente{} \\
    RF-05 & Seleccionar el modo de decisión & \pendiente{prioridad} & El usuario alterna entre el modo equilibrado ($t = 0{,}20$) y el de alta sensibilidad ($t = 0{,}11$), actualizando las decisiones. & \pendiente{} \\
    RF-06 & Exportar resultados del análisis & \pendiente{prioridad} & El sistema genera archivos descargables en formatos CSV y JSON y habilita la vista de impresión. & \pendiente{} \\
    \bottomrule
  \end{tabularx}
\end{table}

Por su parte, la Tabla~\ref{tab:requisitos-no-funcionales} especifica los cinco requisitos no funcionales (RNF-01 a RNF-05) con sus criterios de aceptación y la verificación de la Sección~6.6. En ella constan las pruebas de paridad sobre 120 imágenes en Chromium 158 y Firefox 159 y los tiempos por imagen en un Intel Core i5-12400F (16\,GB RAM). También registra la variación de memoria con 300 imágenes consecutivas y la verificación automática con \texttt{axe-core}~\cite{w3c2023wcag22}.

\begin{table}[H]
  \centering
  \caption{Requisitos no funcionales del aplicativo web, criterios de aceptación y verificación.}
  \label{tab:requisitos-no-funcionales}
  \footnotesize
  \begin{tabularx}{\textwidth}{l >{\hsize=0.75\hsize\linewidth=\hsize}Y >{\hsize=1.10\hsize\linewidth=\hsize}Y >{\hsize=1.15\hsize\linewidth=\hsize}Y}
    \toprule
    ID & Categoría & Criterio de aceptación & Verificación \\
    \midrule
    RNF-01 & Paridad numérica & Píxeles preprocesados y decisiones finales coincidentes con la cadena ONNX de Python en 120 imágenes ($\ge 10$ por clase). & 100\,\% de píxeles idénticos; diferencia máxima en $p_{\text{cal}}(\texttt{mel})$ de $6{,}6 \cdot 10^{-6}$; 100\,\% de coincidencia de clase y decisión en ambos modos en Chromium 158 y Firefox 159. \\
    RNF-02 & Rendimiento & Latencia media de procesamiento e inferencia en CPU local compatible con análisis interactivo y por lotes. & En Intel Core i5-12400F (16\,GB RAM, WebAssembly 4 hilos): carga inicial 1,5\,s y media de 59\,ms/imagen (mediana 52\,ms, p95 79\,ms) en Chromium 158; carga 4,2\,s y 63\,ms/imagen en Firefox 159. \\
    RNF-03 & Estabilidad de memoria & Procesamiento continuo de 300 imágenes seguidas sin fugas de memoria que degraden el navegador. & En Chromium 158 con 300 imágenes consecutivas, la memoria pasó de 96,4\,MB a 99,3\,MB ($+3{,}0$\,\%). \\
    RNF-04 & Accesibilidad & 0 violaciones automáticas en reglas de las Pautas de Accesibilidad para el Contenido Web (WCAG) 2.2 niveles A y AA~\cite{w3c2023wcag22}. & Verificación automática con \texttt{axe-core} (WCAG 2.0, 2.1 y 2.2 niveles A y AA) en 6 estados de la interfaz (incluida vista móvil de 390\,px): 0 violaciones. \\
    RNF-05 & Privacidad & Ejecución local en el cliente sin transmisión de imágenes a servidores externos~\cite{ning2024onnxweb}. & Sitio estático en Vercel; preprocesamiento en OpenCV.js 5.0.0 e inferencia en \texttt{onnxruntime-web} 1.31.1 ejecutados en un \textit{Web Worker} local. \\
    \bottomrule
  \end{tabularx}
\end{table}

\section{Diseño de la interfaz}

El diseño visual e interactivo del aplicativo siguió los criterios de perceptibilidad, operabilidad y comprensibilidad establecidos en la recomendación WCAG 2.2~\cite{w3c2023wcag22}. Los textos e indicadores cumplen el contraste de los niveles A y AA, los controles muestran foco visible y toda la interfaz admite navegación por teclado~\cite{w3c2023wcag22}. Adicionalmente, las probabilidades calibradas y la clase asignada se acompañan de etiquetas descriptivas en español (Tabla~\ref{tab:clases}) para facilitar su lectura.

Según las delimitaciones del estudio~\cite{amaral2025esmo}, la interfaz muestra un aviso visible indicando que el aplicativo es una herramienta de apoyo sobre imágenes dermatoscópicas. Dicho aviso explicita que el sistema no emite diagnósticos médicos ni sustituye la confirmación histopatológica~\cite{amaral2025esmo}. Las Figuras~\ref{fig:vista-principal},~\ref{fig:resultados-lote} y~\ref{fig:detalle-imagen} reservan el espacio para las capturas de la vista principal, la tabla de resultados por lote y el detalle por imagen.

\begin{figure}[H]
  \centering
  \fbox{\parbox[c][3cm][c]{0.85\linewidth}{\centering \pendiente{figura}}}
  \caption{Captura de la vista principal de carga y selección de modo del aplicativo.}
  \label{fig:vista-principal}
\end{figure}

\begin{figure}[H]
  \centering
  \fbox{\parbox[c][3cm][c]{0.85\linewidth}{\centering \pendiente{figura}}}
  \caption{Captura de la tabla de resultados tras el procesamiento de un lote de imágenes.}
  \label{fig:resultados-lote}
\end{figure}

\begin{figure}[H]
  \centering
  \fbox{\parbox[c][3cm][c]{0.85\linewidth}{\centering \pendiente{figura}}}
  \caption{Captura de la vista de detalle de una imagen analizada.}
  \label{fig:detalle-imagen}
\end{figure}

\section{Funcionalidades principales}

\subsection{Carga versátil de imágenes}

El módulo de entrada admite las cuatro modalidades de los requisitos RF-01 a RF-04 (Tabla~\ref{tab:requisitos-funcionales}): archivo individual, archivos múltiples, carpeta local o archivo ZIP. Cuando el lote contiene archivos corruptos o con extensiones distintas de imagen, el aplicativo registra el rechazo con su motivo sin interrumpir el procesamiento de las fotografías válidas restantes. La recepción y validación de los archivos se ejecuta localmente en el navegador sin enviar datos por la red~\cite{ning2024onnxweb}.

\subsection{Modos de decisión configurables}

Para responder al requisito RF-05, la interfaz ofrece un selector que alterna entre el modo equilibrado ($t = 0{,}20$) y el modo de alta sensibilidad ($t = 0{,}11$). Ambos modos aplican la Ecuación~\ref{eq:regla-umbral} sobre las probabilidades calibradas ($T = 3{,}2933$)~\cite{guo2017calibration} para reducir los falsos negativos de melanoma a cambio de más falsos positivos~\cite{barata2023reinforcement,marchetti2023proveai}. El cambio de modo recalcula las decisiones sobre las probabilidades ya obtenidas sin repetir la inferencia del modelo ONNX.

\subsection{Tabla de resultados y vista de detalle}

Al finalizar cada imagen (59\,ms en promedio en Chromium 158), la tabla muestra el archivo, la clase predicha, la alerta de melanoma y $p_{\text{cal}}(\texttt{mel})$. Al seleccionar una fila, se abre la vista de detalle con la miniatura, las probabilidades calibradas de las siete clases de HAM10000~\cite{tschandl2018ham10000} y el umbral activo. Tanto la vista de tabla como la de detalle superaron la verificación automática de accesibilidad en escritorio y en resolución móvil de 390\,px~\cite{w3c2023wcag22}.

\subsection{Exportación de reportes}

Según el requisito RF-06, los resultados pueden descargarse en formato CSV o en formato JSON con las probabilidades de las siete clases y el modo activo. Adicionalmente, el aplicativo dispone de una vista formateada para impresión directa o guardado en documento PDF desde el navegador. Todos los archivos de exportación se generan en la memoria del cliente e incluyen la nota de alcance como herramienta de apoyo~\cite{amaral2025esmo}.

\section{Manual de usuario}

Esta sección describe el procedimiento operativo paso a paso para analizar imágenes dermatoscópicas dentro del aplicativo web. Las instrucciones cubren desde el ingreso a la dirección pública hasta la descarga de reportes en los tres formatos previstos en el requisito RF-06 (Tabla~\ref{tab:requisitos-funcionales}). Para ejecutar una sesión de análisis y exportar los resultados, el usuario debe realizar los cinco pasos numerados que se detallan a continuación:

\begin{enumerate}
  \item \textbf{Acceso al aplicativo:} abrir un navegador compatible con WebAssembly (como Chromium o Firefox) e ingresar a la dirección web \pendiente{URL}, esperando la carga inicial del modelo ONNX (1,5\,s en Chromium 158 o 4,2\,s en Firefox 159).
  \item \textbf{Selección del modo de operación:} elegir en el panel de configuración el modo \textit{Equilibrado} ($t = 0{,}20$, que obtuvo 72,3\,\% de sensibilidad y 89,6\,\% de especificidad para melanoma en la prueba de HAM10000) o el modo \textit{Alta sensibilidad} ($t = 0{,}11$, que obtuvo 80,5\,\% de sensibilidad y 82,3\,\% de especificidad).
  \item \textbf{Carga de imágenes:} seleccionar una de las cuatro opciones de entrada en el área de carga: (a) archivo individual, (b) selección múltiple de archivos, (c) carpeta local o (d) archivo comprimido ZIP; si existen archivos no válidos, revisar el aviso de rechazos generado por la interfaz.
  \item \textbf{Revisión de la tabla y del detalle por imagen:} consultar en la tabla de resultados la clase asignada entre las siete categorías (\texttt{akiec}, \texttt{bcc}, \texttt{bkl}, \texttt{df}, \texttt{mel}, \texttt{nv}, \texttt{vasc}) y pulsar sobre cualquier registro para abrir su vista de detalle con las probabilidades calibradas.
  \item \textbf{Exportación de resultados:} pulsar el botón \textit{Exportar CSV} para descargar la tabla en texto plano, \textit{Exportar JSON} para obtener el registro estructurado de probabilidades o \textit{Imprimir} para generar el reporte imprimible de la sesión.
\end{enumerate}

\section{Verificación del aplicativo}

El protocolo de verificación siguió las recomendaciones de Wang et al.~\cite{wang2025browser}, quienes evaluaron nueve modelos profundos en navegadores de 50 computadores y 20 teléfonos móviles. En ese estudio, la inferencia en navegador fue 16,9 veces más lenta que en CPU nativa y aumentó 67,2\,\% el tiempo de renderizado de la interfaz~\cite{wang2025browser}. Por ello, dichos autores recomiendan medir por separado la carga inicial, la latencia de inferencia, el uso de memoria y la fluidez~\cite{wang2025browser}. Siguiendo esta estructura metodológica, se evaluaron la paridad numérica, la latencia de carga e inferencia, la estabilidad de memoria y la accesibilidad del aplicativo.

La prueba de paridad con Python se ejecutó sobre un subconjunto de 120 imágenes de entrenamiento que incluyó al menos 10 imágenes por cada una de las siete clases diagnósticas. La referencia fue la misma cadena con \texttt{onnxruntime} en Python (no el modelo en Keras), medida en Chromium 158 y verificada con igual resultado en Firefox 159. Las cuatro etapas de preprocesamiento en OpenCV.js 5.0.0 produjeron exactamente los mismos píxeles que en Python, con el 100\,\% de los elementos idénticos. La probabilidad calibrada de melanoma difirió como máximo en $6{,}6 \cdot 10^{-6}$, y la clase final y la decisión de melanoma coincidieron en el 100\,\% de las 120 imágenes en ambos modos.

Las pruebas de rendimiento y memoria se realizaron en un equipo con Intel Core i5-12400F y 16\,GB de RAM, usando \texttt{onnxruntime-web} 1.31.1 en CPU con cuatro hilos~\cite{ning2024onnxweb}. En Chromium 158, la carga inicial tomó 1,5\,s y el procesamiento promedió 59\,ms por imagen (mediana de 52\,ms y percentil 95 de 79\,ms). En Firefox 159, la carga inicial tomó 4,2\,s y el tiempo promedio se situó en 63\,ms por imagen. Con 300 imágenes seguidas en Chromium 158, la memoria pasó de 96,4\,MB a 99,3\,MB ($+3{,}0$\,\%), confirmando la liberación de tensores tras cada imagen.

Finalmente, la accesibilidad se verificó de forma automática con \texttt{axe-core} bajo las reglas WCAG 2.0, 2.1 y 2.2 niveles A y AA~\cite{w3c2023wcag22}. Se evaluaron seis estados de la interfaz: inicio, carga con rechazos, resultados, detalle, acerca del modelo y resultados en móvil de 390\,px. En los seis estados evaluados se registraron 0 violaciones automáticas. De acuerdo con el W3C~\cite{w3c2023wcag22}, esta prueba corresponde a una verificación automática y no acredita la conformidad total con WCAG 2.2 AA sin auditoría manual.

\chapter{Resultados y validación}

Este capítulo presenta las métricas en el conjunto de prueba de HAM10000 para las ocho configuraciones entrenadas y el análisis detallado del modelo seleccionado (MobileNetV3Large a 384\,px, semilla 44). También reporta la validación en las 1.512 imágenes del conjunto de prueba de ISIC 2018~\cite{codella2019isic2018} bajo los tres modos de decisión sin reajustar los umbrales. Por último, desarrolla la discusión de los hallazgos frente a la literatura~\cite{maierhein2024metrics,barata2023reinforcement} y enumera las limitaciones metodológicas, clínicas y técnicas del trabajo.

\section{Resultados de todos los modelos en prueba}

La Tabla~\ref{tab:resultados-modelos} resume el desempeño en la prueba de HAM10000 (1.431 imágenes) para las ocho configuraciones evaluadas con tres semillas (42, 43 y 44) bajo la clase más probable (\textit{argmax}). Como se indicó en la Sección~5.6 (Tabla~\ref{tab:cribado-384}), ConvNeXt-Tiny~\cite{liu2022convnext} se descartó en el cribado a 384\,px (Macro-F1 de 0,729) y no se evaluó en prueba. Por su parte, la Figura~\ref{fig:macrof1-modelos} reserva el espacio para comparar el Macro-F1 medio y la desviación estándar de las ocho configuraciones de la Tabla~\ref{tab:resultados-modelos}.

\begin{sidewaystable}[p]
  \centering
  \caption{Resultados en el conjunto de prueba de HAM10000 (media de tres semillas, clase más probable).}
  \label{tab:resultados-modelos}
  \begin{tabularx}{\linewidth}{>{\hsize=1.50\hsize\linewidth=\hsize}Y >{\hsize=0.90\hsize\linewidth=\hsize}Y >{\hsize=0.75\hsize\linewidth=\hsize}Y >{\hsize=0.95\hsize\linewidth=\hsize}Y >{\hsize=1.10\hsize\linewidth=\hsize}Y >{\hsize=1.00\hsize\linewidth=\hsize}Y >{\hsize=1.05\hsize\linewidth=\hsize}Y >{\hsize=0.75\hsize\linewidth=\hsize}Y}
    \toprule
    Modelo & Resolución (px) & Accuracy & Balanced accuracy & Macro-F1 $\pm$ desv. & Sensibilidad MEL & Especificidad MEL & AUC MEL \\
    \midrule
    MobileNetV3Large & 384 & 82,2\,\% & 62,7\,\% & $0{,}658 \pm 0{,}009$ & 51,4\,\% & 95,9\,\% & 0,900 \\
    EfficientNetV2-S & 384 & 82,0\,\% & 62,2\,\% & $0{,}649 \pm 0{,}007$ & 56,6\,\% & 94,7\,\% & 0,906 \\
    DenseNet121 & 384 & 81,8\,\% & 62,2\,\% & $0{,}649 \pm 0{,}009$ & 54,7\,\% & 94,5\,\% & 0,908 \\
    ResNet50 & 224 & 79,8\,\% & 60,0\,\% & $0{,}632 \pm 0{,}018$ & 44,0\,\% & 95,8\,\% & 0,889 \\
    DenseNet121 & 224 & 79,5\,\% & 61,3\,\% & $0{,}629 \pm 0{,}007$ & 48,2\,\% & 94,4\,\% & 0,883 \\
    DenseNet121-BiLSTM & 224 & 81,0\,\% & 57,5\,\% & $0{,}623 \pm 0{,}024$ & 50,1\,\% & 94,9\,\% & 0,898 \\
    MobileNetV3Large & 224 & 79,1\,\% & 60,0\,\% & $0{,}618 \pm 0{,}005$ & 44,9\,\% & 95,4\,\% & 0,881 \\
    CNN-BiLSTM (desde cero) & 224 & 72,6\,\% & 48,3\,\% & $0{,}492 \pm 0{,}017$ & 38,2\,\% & 94,2\,\% & 0,858 \\
    \bottomrule
  \end{tabularx}
\end{sidewaystable}
\clearpage

\begin{figure}[H]
  \centering
  \fbox{\parbox[c][3cm][c]{0.85\linewidth}{\centering \pendiente{figura}}}
  \caption{Comparación del Macro-F1 medio y desviación estándar entre semillas por configuración en prueba.}
  \label{fig:macrof1-modelos}
\end{figure}

Los datos de la Tabla~\ref{tab:resultados-modelos} muestran que el incremento de resolución de 224\,px a 384\,px mejoró todas las métricas en las arquitecturas evaluadas en ambas resoluciones. En MobileNetV3Large~\cite{howard2019mobilenetv3}, el Macro-F1 medio en prueba subió de $0{,}618 \pm 0{,}005$ a $0{,}658 \pm 0{,}009$ y el AUC de melanoma pasó de 0,881 a 0,900. De igual modo, DenseNet121~\cite{huang2017densenet} incrementó su Macro-F1 medio de $0{,}629 \pm 0{,}007$ a $0{,}649 \pm 0{,}009$ y su AUC de melanoma de 0,883 a 0,908.

Entre las tres arquitecturas entrenadas a 384\,px (MobileNetV3Large con 0,658, EfficientNetV2-S con 0,649 y DenseNet121 con 0,649), la diferencia máxima en Macro-F1 medio fue de 0,009. Esta distancia es inferior a 0,01 y se sitúa en el mismo orden de magnitud que la desviación estándar entre semillas ($\pm 0{,}007$ a $\pm 0{,}009$). Al contrastar la validación (Secciones~5.6 y~5.7) con la prueba de la Tabla~\ref{tab:resultados-modelos}, todos los modelos registraron una caída de 0,07 a 0,10 puntos de Macro-F1. En el extremo inferior de la tabla, la arquitectura CNN-BiLSTM entrenada desde cero obtuvo $0{,}492 \pm 0{,}017$ de Macro-F1, 48,3\,\% de balanced accuracy y 38,2\,\% de sensibilidad para melanoma.

\section{Modelo seleccionado}

El modelo elegido por la regla de selección preregistrada en validación (Sección~5.7) fue MobileNetV3Large a 384\,px con la semilla 44~\cite{howard2019mobilenetv3}. En las 1.431 imágenes de prueba de HAM10000 (clase más probable), obtuvo 81,8\,\% de accuracy, 61,7\,\% de balanced accuracy y 64,7\,\% de Macro-F1 (IC 95\,\% por \textit{bootstrap}: 58,5--69,5\,\%). En la discriminación binaria de melanoma frente al resto de clases, el clasificador alcanzó un AUC de 0,903 (IC 95\,\% por \textit{bootstrap}: 0,879--0,924).

La Tabla~\ref{tab:modos-decision} compara el modelo seleccionado en la prueba de HAM10000 bajo la clase más probable y los dos modos calibrados ($T = 3{,}2933$). Al pasar al modo equilibrado ($t = 0{,}20$), la sensibilidad para melanoma subió de 50,9\,\% a 72,3\,\% ($+21{,}4$ puntos), con 89,6\,\% de especificidad y 64,1\,\% de Macro-F1. En el modo de alta sensibilidad ($t = 0{,}11$), la sensibilidad para melanoma se elevó hasta 80,5\,\% a cambio de una especificidad de 82,3\,\% y una accuracy global de 75,5\,\%.

\begin{table}[H]
  \centering
  \caption{Desempeño del modelo seleccionado (MobileNetV3Large, 384\,px, semilla 44) por modo de decisión en la prueba de HAM10000.}
  \label{tab:modos-decision}
  \begin{tabularx}{\textwidth}{>{\hsize=1.35\hsize\linewidth=\hsize}Y >{\hsize=0.80\hsize\linewidth=\hsize}Y >{\hsize=1.25\hsize\linewidth=\hsize}Y >{\hsize=0.80\hsize\linewidth=\hsize}Y >{\hsize=0.80\hsize\linewidth=\hsize}Y}
    \toprule
    Modo de decisión & Accuracy & Macro-F1 & Sensibilidad MEL & Especificidad MEL \\
    \midrule
    Clase más probable & 81,8\,\% & 64,7\,\% (IC 95\,\%: 58,5--69,5) & 50,9\,\% & 95,5\,\% \\
    Equilibrado ($t = 0{,}20$) & 79,9\,\% & 64,1\,\% & 72,3\,\% & 89,6\,\% \\
    Alta sensibilidad ($t = 0{,}11$) & 75,5\,\% & 62,0\,\% & 80,5\,\% & 82,3\,\% \\
    \bottomrule
  \end{tabularx}
\end{table}

Para examinar el comportamiento desagregado entre las siete categorías diagnósticas~\cite{maierhein2024metrics}, la Tabla~\ref{tab:f1-clase} reporta el puntaje F1 por clase del modelo seleccionado bajo la regla de la clase más probable. La clase mayoritaria \texttt{nv} alcanzó un F1 de 0,92 y la clase minoritaria \texttt{vasc} obtuvo el segundo valor más alto con 0,74. En cambio, los menores puntajes F1 se registraron en \texttt{akiec} (0,55), \texttt{mel} (0,55) y \texttt{df} (0,48).

\begin{table}[H]
  \centering
  \caption{Puntaje F1 por clase del modelo seleccionado (MobileNetV3Large, 384\,px, semilla 44, clase más probable).}
  \label{tab:f1-clase}
  \begin{tabularx}{\textwidth}{Yr}
    \toprule
    Clase & F1 \\
    \midrule
    \texttt{nv} & 0,92 \\
    \texttt{vasc} & 0,74 \\
    \texttt{bkl} & 0,66 \\
    \texttt{bcc} & 0,63 \\
    \texttt{akiec} & 0,55 \\
    \texttt{mel} & 0,55 \\
    \texttt{df} & 0,48 \\
    \bottomrule
  \end{tabularx}
\end{table}

La Figura~\ref{fig:matriz-confusion} reserva el espacio para la matriz de confusión del modelo seleccionado en el conjunto de prueba de HAM10000 (archivos \texttt{docs/resultados/fase\_e2/deployable\_test\_cm\_*.csv}). Con la clase más probable, acertó 889/958 \texttt{nv}, 112/157 \texttt{bkl}, 45/74 \texttt{bcc}, 23/46 \texttt{akiec}, 13/20 \texttt{vasc} y 7/17 \texttt{df}. En la clase \texttt{mel}, el modelo acertó 81 de los 159 melanomas de prueba; de los 78 falsos negativos restantes, 50 se predijeron como \texttt{nv} y 21 como \texttt{bkl}. En el modo equilibrado ($t = 0{,}20$) detectó 115 de 159 melanomas (con 87 \texttt{nv} marcados como \texttt{mel}), y en alta sensibilidad ($t = 0{,}11$) detectó 128 de 159 (con 147 \texttt{nv}). Por su parte, la Figura~\ref{fig:roc-melanoma} reserva el espacio para la curva ROC de melanoma (\texttt{docs/resultados/fase\_e2/roc\_mel.png}), sobre la cual se obtuvo el AUC de 0,903.

\begin{figure}[H]
  \centering
  \fbox{\parbox[c][3cm][c]{0.85\linewidth}{\centering \pendiente{figura}}}
  \caption{Matriz de confusión del modelo seleccionado en el conjunto de prueba de HAM10000.}
  \label{fig:matriz-confusion}
\end{figure}

\begin{figure}[H]
  \centering
  \fbox{\parbox[c][3cm][c]{0.85\linewidth}{\centering \pendiente{figura}}}
  \caption{Curva ROC para la clase melanoma del modelo seleccionado en la prueba de HAM10000 (AUC $= 0{,}903$).}
  \label{fig:roc-melanoma}
\end{figure}

\section{Validación en ISIC 2018}

Para evaluar la generalización fuera de la partición de HAM10000, el modelo seleccionado (MobileNetV3Large, 384\,px, semilla 44) se probó en las 1.512 imágenes de prueba de ISIC 2018 Task 3~\cite{codella2019isic2018}. Con los identificadores del ISIC Archive, se comprobó que ninguna de las 1.512 imágenes compartiera lesión con las 7.153 imágenes de entrenamiento. La evaluación mantuvo fijos y sin reajuste tanto el parámetro de temperatura $T = 3{,}2933$ como los umbrales de decisión $t = 0{,}20$ y $t = 0{,}11$. La prueba en HAM10000 se ejecutó con el modelo en Keras y la de ISIC 2018 con la cadena del aplicativo (\pendiente{confirmar}).

La Tabla~\ref{tab:validacion-isic} contrasta las métricas obtenidas en la prueba de HAM10000 (1.431 imágenes) y en la prueba de ISIC 2018 (1.512 imágenes) para los tres modos de operación. En ISIC 2018, el AUC de melanoma fue de 0,908 (0,903 en HAM10000), el Macro-F1 subió a 67,9\,\% (frente a 64,7\,\%) y la balanced accuracy pasó de 61,7\,\% a 64,8\,\%. La sensibilidad para melanoma en ISIC 2018 fue de 53,8\,\% con la clase más probable, de 74,9\,\% en modo equilibrado (especificidad 87,7\,\%) y de 87,1\,\% en alta sensibilidad (especificidad 79,3\,\%).

\begin{sidewaystable}[p]
  \centering
  \caption{Comparación del modelo seleccionado en el conjunto de prueba de HAM10000 (1.431 imágenes) y en ISIC 2018 Task 3 Test (1.512 imágenes).}
  \label{tab:validacion-isic}
  \begin{tabularx}{\linewidth}{>{\hsize=1.05\hsize\linewidth=\hsize}Y >{\hsize=1.20\hsize\linewidth=\hsize}Y >{\hsize=0.85\hsize\linewidth=\hsize}Y >{\hsize=0.95\hsize\linewidth=\hsize}Y >{\hsize=0.90\hsize\linewidth=\hsize}Y >{\hsize=1.05\hsize\linewidth=\hsize}Y >{\hsize=1.10\hsize\linewidth=\hsize}Y >{\hsize=0.90\hsize\linewidth=\hsize}Y}
    \toprule
    Conjunto & Modo & Accuracy & Balanced accuracy & Macro-F1 & Sensibilidad MEL & Especificidad MEL & AUC MEL \\
    \midrule
    HAM10000 & Clase más probable & 81,8\,\% & 61,7\,\% & 64,7\,\% & 50,9\,\% & 95,5\,\% & 0,903 \\
    ISIC 2018 & Clase más probable & 78,8\,\% & 64,8\,\% & 67,9\,\% & 53,8\,\% & \pendiente{} & 0,908 \\
    HAM10000 & Equilibrado ($t = 0{,}20$) & 79,9\,\% & 61,8\,\% & 64,1\,\% & 72,3\,\% & 89,6\,\% & 0,903 \\
    ISIC 2018 & Equilibrado ($t = 0{,}20$) & 75,6\,\% & 64,6\,\% & 66,4\,\% & 74,9\,\% & 87,7\,\% & 0,908 \\
    HAM10000 & Alta sensibilidad ($t = 0{,}11$) & 75,5\,\% & 60,0\,\% & 62,0\,\% & 80,5\,\% & 82,3\,\% & 0,903 \\
    ISIC 2018 & Alta sensibilidad ($t = 0{,}11$) & 70,8\,\% & 62,5\,\% & 63,7\,\% & 87,1\,\% & 79,3\,\% & 0,908 \\
    \bottomrule
  \end{tabularx}
\end{sidewaystable}
\clearpage

La caída de accuracy en ISIC 2018 frente a HAM10000 (de 81,8\,\% a 78,8\,\% con la clase más probable) obedece en parte a la diferente composición de clases~\cite{codella2019isic2018,maierhein2024metrics}. En HAM10000 la clase \texttt{nv} representa el 67,0\,\% de las imágenes, mientras que en ISIC 2018 constituye cerca del 60\,\%, diferencia que explica unos 2 puntos de caída en accuracy. El mejor resultado del reto ISIC 2018 alcanzó una balanced accuracy de 0,885~\cite{codella2019isic2018}, no comparable directamente con 0,648 (Tabla~\ref{tab:validacion-isic}) porque empleó ensambles y datos externos adicionales a HAM10000.

Frente a Barata et al.~\cite{barata2023reinforcement} (cuya sensibilidad para melanoma en ISIC 2018 subió de 61,4\,\% a 79,5\,\%), los modos calibrados pasaron de 53,8\,\% a 74,9\,\% y 87,1\,\%. Al comparar debe considerarse que Barata et al.~\cite{barata2023reinforcement} evaluaron 1.511 imágenes, mientras que este trabajo evaluó las 1.512 del conjunto (\pendiente{indicar si la imagen ISIC\_0035068 está entre las 1.512 de ISIC 2018 y cómo se trató}). Además, los metadatos públicos no separan las imágenes de ISIC 2018 provenientes de centros nuevos de las capturadas en los mismos centros de HAM10000~\cite{tschandl2018ham10000,codella2019isic2018}.

\section{Discusión}

El primer aspecto observado en el modelo seleccionado (Tabla~\ref{tab:modos-decision}) es la brecha de 20,1 puntos porcentuales entre la accuracy (81,8\,\%) y la balanced accuracy (61,7\,\%) con la clase más probable. Esta diferencia proviene del desbalance de HAM10000, donde \texttt{nv} concentra el 67,0\,\% de la prueba (958 de 1.431 imágenes) y alcanza un F1 de 0,92 (Tabla~\ref{tab:f1-clase}). Como señalan Maier-Hein et al.~\cite{maierhein2024metrics}, una clase mayoritaria dominante hace que la accuracy refleje el acierto sobre ella y oculte los fallos en clases pequeñas, justificando el uso de Macro-F1.

El segundo punto concierne al patrón de error en la clase \texttt{mel} bajo la regla estándar de la clase más probable. De los 159 melanomas de prueba de HAM10000, el modelo acertó 81 (50,9\,\% de sensibilidad) y tuvo 78 fallos, de los cuales 50 fueron a \texttt{nv} y 21 a \texttt{bkl}. Sin embargo, sobre esas mismas 1.431 imágenes el clasificador obtuvo un AUC de melanoma de 0,903 (IC 95\,\%: 0,879--0,924) y una especificidad de 95,5\,\%. Tener un AUC de 0,903 con 50,9\,\% de sensibilidad indica un problema de punto de corte: \textit{argmax} exige que $p(\texttt{mel})$ supere a \texttt{nv} y \texttt{bkl}, penalizando las lesiones limítrofes.

El tercer elemento de análisis es el efecto práctico de los umbrales calibrados ($t = 0{,}20$ y $t = 0{,}11$) sobre el balance entre falsos negativos y falsos positivos de melanoma~\cite{barata2023reinforcement}. Con $t = 0{,}20$ (modo equilibrado), los melanomas detectados en HAM10000 subieron de 81 a 115 de 159 (72,3\,\%), y los nevos marcados como melanoma pasaron de 38 a 87 (especificidad 89,6\,\%). Con $t = 0{,}11$ (alta sensibilidad), la detección subió a 128 de 159 melanomas (80,5\,\%) y los nevos marcados como melanoma pasaron a 147 (especificidad 82,3\,\%). Este compromiso coincide con estudios prospectivos como PROVE-AI (96,8\,\% de sensibilidad y 37,4\,\% de especificidad)~\cite{marchetti2023proveai} y Heinlein et al. (92,1\,\% y 67,3\,\%)~\cite{heinlein2024multicenter}, donde mayor sensibilidad implica menor especificidad.

En cuarto lugar, el desglose por categoría de la Tabla~\ref{tab:f1-clase} muestra que el tamaño de muestra en entrenamiento no es el único determinante del puntaje F1 por clase. Con 102 imágenes de entrenamiento, \texttt{vasc} logró el segundo mayor F1 en prueba (0,74, acertando 13 de 20); como hipótesis, una explicación posible es su color rojizo o violáceo distintivo~\cite{tschandl2018ham10000}. En cambio, \texttt{df} tuvo el menor F1 (0,48, 7 de 17 aciertos) tras entrenarse con 82 imágenes. A su lado, \texttt{akiec} obtuvo un F1 de 0,55 (23 de 46 aciertos) al confundirse en 10 casos con \texttt{bkl}~\cite{tschandl2018ham10000}.

El quinto punto contrasta el Macro-F1 de 0,647 (accuracy de 81,8\,\%) de este trabajo con los valores mayores al 90\,\% de accuracy reportados en parte de la literatura sobre HAM10000~\cite{hu2024multiscale}. En numerosos estudios previos no se verifica una partición agrupada por \texttt{lesion\_id} o se aplica sobremuestreo antes de dividir los subconjuntos~\cite{abhishek2025quality,kapoor2023leakage}. En la Tabla~\ref{tab:comparacion-particiones}, partir por imagen dejó el 36,1\,\% de la prueba y el 63,5\,\% de los melanomas en entrenamiento, subiendo el Macro-F1 de $+8{,}7$ a $+9{,}2$ puntos. Sin señalar errores en trabajos concretos sin código público, estos datos muestran la dirección de la inflación producida por compartir lesiones entre particiones~\cite{kapoor2023leakage,abhishek2025quality}.

Finalmente, el sexto aspecto se refiere al alcance de la calibración mediante \textit{temperature scaling}~\cite{guo2017calibration,sambyal2023calibration}. El ajuste de $T = 3{,}2933$ en validación suavizó las probabilidades de salida y permitió trasladar sin cambios los umbrales $t = 0{,}20$ y $t = 0{,}11$ al conjunto ISIC 2018 (Tabla~\ref{tab:validacion-isic}). No obstante, según Sambyal et al.~\cite{sambyal2023calibration} y Penso et al.~\cite{penso2024calibration}, escalar por un único parámetro $T$ no altera las confusiones entre clases cercanas y depende de las etiquetas. Tampoco convierte la probabilidad calibrada en una estimación de riesgo clínico individualizado~\cite{sambyal2023calibration,penso2024calibration}.

\section{Limitaciones}

El análisis de los resultados experimentales y del prototipo web debe interpretarse dentro de diez limitaciones metodológicas, de datos y de validación clínica~\cite{collins2024tripodai,vasey2022decideai}. Cada una de estas restricciones acota las conclusiones del estudio y define los requerimientos para futuras iteraciones del sistema. A continuación se enumeran las limitaciones identificadas junto con su consecuencia directa:

\begin{itemize}
  \item \textbf{Evaluación sobre una única partición agrupada:} los experimentos se ejecutaron sobre una sola división fija 5/7--1/7--1/7 (semilla 42) en lugar de un esquema de validación cruzada de $k$ pliegues agrupados, lo que impide estimar la varianza del rendimiento frente a distintas asignaciones de lesiones entre entrenamiento y prueba.
  \item \textbf{Doble apertura del conjunto de prueba:} el conjunto de prueba de HAM10000 se utilizó en dos momentos del proyecto (al concluir la etapa de 224\,px y al concluir la etapa de 384\,px); aunque en ambas ocasiones la regla de selección estaba fijada en validación antes de abrirlo, el reuso del conjunto de prueba reduce la independencia estricta de la segunda evaluación~\cite{kapoor2023leakage}.
  \item \textbf{Tamaño reducido de las clases minoritarias en prueba:} las categorías \texttt{df}, \texttt{vasc} y \texttt{akiec} cuentan únicamente con 17, 20 y 46 imágenes en el conjunto de prueba (Tabla~\ref{tab:particion}), lo que amplía la incertidumbre estadística de sus métricas individuales y explica parte del ancho del intervalo de confianza del Macro-F1 (58,5\,\%--69,5\,\%).
  \item \textbf{Meta inicial de Macro-F1 no alcanzada:} el modelo seleccionado obtuvo un Macro-F1 en prueba de 0,647, valor inferior a la meta de 0,80 planteada inicialmente en el anteproyecto, como consecuencia de evaluar bajo una partición estricta por lesión sin fuga de datos y con siete clases desbalanceadas (Tabla~\ref{tab:distribucion}).
  \item \textbf{Restricción a imágenes dermatoscópicas sin metadatos clínicos:} el clasificador opera exclusivamente con píxeles de imágenes dermatoscópicas y no incorpora variables clínicas del paciente (como edad, sexo, antecedentes o localización anatómica) ni es aplicable a fotografías macroscópicas tomadas con teléfonos móviles~\cite{amaral2025esmo}.
  \item \textbf{Agrupación por lesión y no por paciente, con duplicados residuales:} dado que HAM10000 no incluye un identificador único de paciente~\cite{tschandl2018ham10000}, agrupar por \texttt{lesion\_id} no impide que lesiones distintas de una misma persona queden en subconjuntos diferentes; además, la auditoría contra los 18 pares de~\cite{abhishek2025quality} detectó 2 imágenes de prueba de la clase \texttt{bkl} (0,14\,\%) que tienen un duplicado con distinto \texttt{lesion\_id} en entrenamiento.
  \item \textbf{Validación externa parcialmente solapada en origen institucional:} aunque se verificó que ninguna de las 1.512 imágenes de ISIC 2018 compartiera lesión con el entrenamiento, dicho conjunto incluye muestras de las mismas instituciones que aportaron a HAM10000 y no fue posible aislar las imágenes de centros nuevos~\cite{codella2019isic2018,tschandl2018ham10000}.
  \item \textbf{Representatividad demográfica y de fototipos cutáneos:} los datos de HAM10000 proceden de poblaciones de Austria y Australia con predominio de pieles claras~\cite{tschandl2018ham10000}, por lo que el desempeño del modelo no puede extrapolarse a fototipos oscuros sin una validación específica que controle los sesgos documentados en la literatura~\cite{daneshjou2022disparities}.
  \item \textbf{Ausencia de validación clínica prospectiva:} todas las pruebas realizadas son retrospectivas en laboratorio (TRL 4) y no evalúan la interacción del médico con la herramienta en la consulta real ni su impacto en las decisiones clínicas según los lineamientos DECIDE-AI~\cite{vasey2022decideai,marchetti2023proveai,heinlein2024multicenter}.
  \item \textbf{Verificación de accesibilidad exclusivamente automática:} la inspección con \texttt{axe-core} cubrió solo las reglas verificables de forma programática en seis estados de la interfaz, por lo que no sustituye una evaluación manual con tecnologías de asistencia para certificar la totalidad de WCAG 2.2 AA~\cite{w3c2023wcag22}.
\end{itemize}

\chapter{Conclusiones}

Este capítulo presenta las conclusiones del trabajo de grado organizadas según los cuatro objetivos específicos planteados en la Sección~2.3. Cada apartado sintetiza las actividades ejecutadas y las cifras experimentales verificadas en los capítulos previos (Capítulos~4 a~7). Adicionalmente, se incluye una conclusión metodológica sobre el efecto de la unidad de partición en la evaluación de modelos dermatoscópicos.

\begin{itemize}
  \item \textbf{Gestión, preprocesamiento y aumento de la base de datos (Objetivo específico 1):} se procesaron las 10.015 imágenes dermatoscópicas correspondientes a 7.470 lesiones de las siete clases del repositorio HAM10000~\cite{tschandl2018ham10000} (Tabla~\ref{tab:distribucion}). Mediante \texttt{StratifiedGroupKFold}~\cite{pedregosa2011sklearn}, los datos se dividieron por \texttt{lesion\_id} en 7.153 imágenes de entrenamiento, 1.431 de validación y 1.431 de prueba sin lesiones compartidas entre subconjuntos (Tabla~\ref{tab:particion}). A todas las imágenes se les aplicó redimensionado con relleno negro, filtro de mediana $3 \times 3$, remoción de vello~\cite{lee1997dullrazor,telea2004inpainting} y ecualización CLAHE con límite 2,0~\cite{zuiderveld1994clahe}. El desbalance se trató en el entrenamiento mediante aumento geométrico y fotométrico en línea junto con pesos de clase entre 0,46 y 3,53 (Tabla~\ref{tab:pesos}).

  \item \textbf{Implementación de arquitecturas basadas en redes neuronales convolucionales (Objetivo específico 2):} se implementaron siete arquitecturas de clasificación de imágenes dermatoscópicas evaluadas en dos resoluciones de entrada de 224\,px y 384\,px (Tabla~\ref{tab:arquitecturas}). El conjunto de modelos incluyó cinco redes preentrenadas en ImageNet~\cite{deng2009imagenet} (DenseNet121~\cite{huang2017densenet}, ResNet50~\cite{he2016resnet}, MobileNetV3Large~\cite{howard2019mobilenetv3}, EfficientNetV2-S~\cite{tan2021efficientnetv2} y ConvNeXt-Tiny~\cite{liu2022convnext}), una red híbrida DenseNet121-BiLSTM y una red CNN-BiLSTM entrenada desde cero. Las redes con transferencia de aprendizaje se ajustaron con una receta de dos fases (hasta 15 épocas con base congelada y hasta 40 épocas de ajuste fino con decaimiento coseno). Ocho configuraciones se entrenaron con tres semillas aleatorias (42, 43 y 44) y se compararon en el conjunto de prueba de 1.431 imágenes (Tabla~\ref{tab:resultados-modelos}).

  \item \textbf{Evaluación y selección del modelo para identificación de melanoma (Objetivo específico 3):} mediante la regla de selección preregistrada en validación ($\delta = 0{,}0112$ y desempate por AUC de melanoma), se eligió MobileNetV3Large a 384\,px con semilla 44~\cite{howard2019mobilenetv3}. En la prueba de HAM10000 con la clase más probable, este modelo obtuvo 81,8\,\% de accuracy, 0,647 de Macro-F1 y 0,903 de AUC de melanoma (Tabla~\ref{tab:modos-decision}). En la validación sobre las 1.512 imágenes de prueba de ISIC 2018~\cite{codella2019isic2018}, el modelo alcanzó un AUC de melanoma de 0,908 y un Macro-F1 de 67,9\,\% sin modificar sus parámetros (Tabla~\ref{tab:validacion-isic}). Con la calibración por temperatura ($T = 3{,}2933$)~\cite{guo2017calibration}, la sensibilidad para melanoma en HAM10000 alcanzó 72,3\,\% en el modo equilibrado ($t = 0{,}20$) y 80,5\,\% en alta sensibilidad ($t = 0{,}11$).

  \item \textbf{Integración del modelo seleccionado en el aplicativo web (Objetivo específico 4):} el modelo seleccionado se exportó al formato ONNX (12,9\,MB) y se integró en un aplicativo web estático desarrollado con React, TypeScript y Vite (Sección~6.1). El aplicativo ejecuta en un Web Worker el preprocesamiento con OpenCV.js 5.0.0 y la inferencia con ONNX Runtime Web 1.31.1~\cite{ning2024onnxweb} sin transmitir imágenes fuera del navegador. En las 120 imágenes de prueba de paridad, el aplicativo coincidió en el 100\,\% de las decisiones con la cadena de referencia en Python con \texttt{onnxruntime} (Sección~6.6). El procesamiento tomó en promedio 59\,ms por imagen en Chromium 158, se verificó con cero violaciones automáticas en \texttt{axe-core} y se desplegó como sitio web estático.

  \item \textbf{Hallazgo metodológico sobre la unidad de partición:} la comparación entre esquemas de división en HAM10000 mostró que partir por imagen deja el 36,1\,\% de la prueba con otra fotografía de la misma lesión en entrenamiento (Sección~5.5). En la clase \texttt{mel}, esta superposición entre entrenamiento y prueba afectó al 63,5\,\% de los melanomas evaluados (106 de 167 imágenes). Como resultado, la partición por imagen infló el Macro-F1 de prueba cerca de 9 puntos porcentuales ($+8{,}7$ puntos en DenseNet121 y $+9{,}2$ puntos en CNN-BiLSTM, Tabla~\ref{tab:comparacion-particiones})~\cite{kapoor2023leakage,abhishek2025quality}.
\end{itemize}

\chapter{Trabajo futuro}

Este capítulo describe las líneas de trabajo futuro derivadas de las limitaciones metodológicas, clínicas y de ingeniería documentadas en la Sección~7.5. Cada propuesta se vincula con una restricción observada durante el desarrollo experimental o en la verificación del prototipo web (Capítulos~4 a~7).

Para superar la evaluación sobre una única partición fija 5/7--1/7--1/7 (Sección~4.3), se propone ejecutar validación cruzada de $k$ pliegues agrupados mediante \texttt{StratifiedGroupKFold}~\cite{pedregosa2011sklearn}. Este esquema cuantificará la varianza del Macro-F1 entre distintas particiones y evitará reutilizar un mismo conjunto de prueba en múltiples etapas experimentales~\cite{kapoor2023leakage}. Complementariamente, el uso de bancos con identificadores de paciente e inspección de duplicados eliminará el solapamiento residual del 0,14\,\% detectado en HAM10000~\cite{abhishek2025quality}.

Frente a la restricción de operar únicamente con píxeles dermatoscópicos (Sección~7.5), una segunda línea consiste en incorporar metadatos clínicos como edad, sexo y localización anatómica~\cite{tschandl2018ham10000}. La fusión multimodal de estas variables con los mapas de características de MobileNetV3Large~\cite{howard2019mobilenetv3} aportaría contexto clínico para desambiguar lesiones melanocíticas y queratósicas~\cite{amaral2025esmo}. En paralelo, un estudio de ablación por etapa medirá el aporte individual del filtro de mediana, la remoción de vello~\cite{wang2024hair} y el ajuste CLAHE~\cite{zuiderveld1994clahe}.

Para atender el reducido tamaño de prueba en \texttt{df} (17 imágenes), \texttt{vasc} (20) y \texttt{akiec} (46) (Tabla~\ref{tab:particion}), se plantea ampliar el entrenamiento con repositorios multicéntricos mayores~\cite{codella2019isic2018}. De igual forma, dado que HAM10000 proviene de poblaciones de Austria y Australia con predominio de pieles claras~\cite{tschandl2018ham10000}, se requiere incorporar imágenes de fototipos oscuros subrepresentados~\cite{daneshjou2022disparities}. Como demostraron Daneshjou et al.~\cite{daneshjou2022disparities}, el ajuste con cohortes diversas permite medir y reducir las brechas de desempeño entre subgrupos demográficos.

Con respecto a la brecha frente a la meta inicial de 0,80 de Macro-F1 (0,647 en la Tabla~\ref{tab:modos-decision}), se propone reemplazar la inicialización en ImageNet~\cite{deng2009imagenet}. El ajuste fino puede partir de modelos fundacionales específicos del dominio como PanDerm, preentrenado con más de dos millones de imágenes dermatológicas de once instituciones~\cite{yan2025panderm}. Esta transferencia desde representaciones dermatológicas previas debe evaluarse junto con técnicas de cuantización para conservar un tamaño compatible con la ejecución en navegador~\cite{ning2024onnxweb}.

En relación con la pérdida de especificidad observada al reducir el umbral escalar de melanoma a $t = 0{,}11$ (Tabla~\ref{tab:modos-decision}), se plantea explorar reglas de decisión sensibles al costo clínico~\cite{barata2023reinforcement}. A diferencia del corte sobre una sola probabilidad calibrada~\cite{guo2017calibration}, el enfoque de aprendizaje por refuerzo de Barata et al.~\cite{barata2023reinforcement} asigna recompensas asimétricas según la severidad de cada confusión diagnóstica. Incorporar una matriz de costos clínicos en las siete clases buscaría reducir los 50 melanomas confundidos con \texttt{nv} y los 21 confundidos con \texttt{bkl} (Sección~7.2).

Dado que el sistema se encuentra en nivel de madurez TRL~4 con pruebas retrospectivas (Sección~2.5), el paso necesario hacia un entorno asistencial es una evaluación clínica inicial supervisada~\cite{vasey2022decideai}. Siguiendo la guía DECIDE-AI~\cite{vasey2022decideai}, dicho estudio debe analizar la interacción de médicos y dermatólogos con los modos equilibrado y de alta sensibilidad en flujos reales de trabajo. Esta validación prospectiva requerirá cohortes de instituciones independientes de HAM10000 e ISIC 2018 y confirmación histopatológica de las lesiones sospechosas~\cite{amaral2025esmo,marchetti2023proveai,heinlein2024multicenter}.

Para formalizar la documentación de futuras versiones del clasificador y su validación en nuevos centros, se propone estructurar los informes técnicos bajo la guía TRIPOD+AI~\cite{collins2024tripodai}. La lista de verificación de TRIPOD+AI establece criterios estandarizados para reportar la procedencia de los datos, el tratamiento de valores faltantes, la calibración~\cite{sambyal2023calibration,penso2024calibration} y los protocolos de evaluación~\cite{collins2024tripodai}. Auditar cada ítem de dicha guía facilitará la reproducibilidad independiente de los experimentos y la comparación entre iteraciones del modelo~\cite{collins2024tripodai,maierhein2024metrics}.

Finalmente, para subsanar el alcance parcial de la verificación automática con \texttt{axe-core} (Sección~6.6), se plantea realizar una auditoría manual de accesibilidad sobre el aplicativo web~\cite{w3c2023wcag22}. Esta evaluación debe incluir pruebas con lectores de pantalla, recorrido por teclado sin ratón y validación con usuarios en distintos dispositivos móviles y de escritorio~\cite{w3c2023wcag22,wang2025browser}. Solo la combinación de pruebas automáticas y verificación humana permite evaluar la totalidad de los criterios de éxito de nivel A y AA de WCAG~2.2~\cite{w3c2023wcag22}.


\clearpage
\phantomsection
\chapter*{Anexo: Declaración sobre el uso de Inteligencia Artificial}
\addcontentsline{toc}{chapter}{Anexo: Declaración sobre el uso de Inteligencia Artificial}
\markboth{}{}
% GUÍA: El autor redactará la declaración a partir de su uso real de herramientas y de sus verificaciones; no presuponer herramientas, tareas ni alcance de autoría.
\section*{Herramientas usadas}
% GUÍA: Identificar herramientas efectivamente usadas. Fuente: registro personal del autor.
\pendiente{herramientas usadas}

\section*{Tareas en que se usaron}
% GUÍA: Describir las tareas concretas asistidas. Fuente: registro de trabajo del autor.
\pendiente{tareas en que se usaron}

\section*{Qué hizo y verificó el autor}
% GUÍA: Precisar elaboración propia y verificaciones realizadas. Fuente: evidencias y revisión del autor.
\pendiente{qué hizo y verificó el autor}

\section*{Nota de autoría}
% GUÍA: Redactar personalmente la nota de autoría. Fuente: declaración del autor.
\pendiente{nota de autoría}


% GUÍA: Referencias recuperadas de la base bibliográfica previa del proyecto y corregidas.
% El PDF disponible es una plantilla sin bibliografía; cotejar con el anteproyecto real.
% Se muestran todas para facilitar la revisión; sustituir nocite por citas en la redacción final.
\bibliographystyle{ieeetr}
\bibliography{Referencias_melanoma_2022_en_adelante,fundamentales,referencias}
\end{document}
